
Taken from \url{https://holgergerhardt.github.io/scbr/2024-10-13_estimation-illustration.R}:

\begin{lstlisting}
# This file illustrates the most important parts of the estimation of
# preference parameters performed by Holger Gerhardt & Rafael Suchy (2024),
# “Estimating Preference Parameters from Strictly Concave Budget Restrictions,”
# https://www.econtribute.de/RePEc/ajk/ajkdps/ECONtribute_336_2024.pdf.

# Version: 2024-10-13
# The results reported in this file were produced using R version 4.4.1
# (https://cloud.r-project.org/bin/macosx/big-sur-arm64/base/R-4.4.1-arm64.pkg)
# and the up-to-date versions of the packages mentioned below.




# PREAMBLE ------------------------------------------------------------------------------------


# Clean up environment
rm(list = ls())

options(max.print = 9999)  # Increase limit for omitting entries in output
options(scipen = 999)  # Increase limit of using scientific notation
round_prec <- 4  # Number of decimal places for reporting the parameter estimates

# Required packages
packages_required <- c(
  "AER",  # For "tobit() function"
  "cli",  # For colored error and warning messages
  "msm",  # For "deltamethod() function"
  "maxLik"  # For maximum likelihood estimation
)
# Install packages that are not installed yet
install.packages(setdiff(packages_required, rownames(installed.packages())))
# For updating already installed packages, use the following
# install.packages(packages_required)
# Load the required packages
for (name in packages_required) {
  library(name, character.only = TRUE)
}

# Alternatively, use groundhog:
# install.packages("groundhog")
# library("groundhog")
# groundhog.library(packages_required, "2024-10-10")




# BUDGET RESTRICTIONS -------------------------------------------------------------------------


# General functional form, based on Equation (9) from
# https://www.econtribute.de/RePEc/ajk/ajkdps/ECONtribute_336_2024.pdf:
# c_{t}^{1 + z} + (1 / R)^{1 + z} c_{t + k}^{1 + z} = m^{1 + z}.
# Thus, linear budget restrictions (LBRs) for z = 0, and
# strictly concave budget restrictions (SCBRs) for z > 0.
c_2 <- function(c_1, m, R, z, b_1, b_2) {
  R * (m^(1 + z) - c_1^(1 + z))^(1 / (1 + z))
}

# Experimental parameters
# (see Table E.1 in https://www.econtribute.de/RePEc/ajk/ajkdps/ECONtribute_336_2024.pdf)
# The presentation of the different budget restrictions in the experiment by Gerhardt & Suchy
# was randomized as described in the manuscript. We abstract from this randomization here.
budget_restritions <- matrix(c(
  c(01, 1.50, 1.50, 0.0, 1, 05, 020.00, 0.70000, 1.42857, 1, 1, 0.01, 02.5, 025),
  c(02, 1.50, 1.50, 0.0, 1, 05, 017.50, 0.80000, 1.25000, 1, 1, 0.01, 02.5, 025),
  c(03, 1.50, 1.50, 0.0, 1, 05, 015.56, 0.90000, 1.11111, 1, 1, 0.01, 02.5, 025),
  c(04, 1.50, 1.50, 0.0, 1, 05, 014.70, 0.95238, 1.05000, 1, 1, 0.01, 02.5, 025),
  c(05, 1.50, 1.50, 0.0, 1, 05, 014.00, 1.00000, 1.00000, 1, 1, 0.01, 02.5, 025),
  c(06, 1.50, 1.50, 0.0, 1, 05, 014.00, 1.05000, 0.95238, 1, 1, 0.01, 02.5, 025),
  c(07, 1.50, 1.50, 0.0, 1, 05, 014.00, 1.11111, 0.90000, 1, 1, 0.01, 02.5, 025),
  c(08, 1.50, 1.50, 0.0, 1, 05, 014.00, 1.25000, 0.80000, 1, 1, 0.01, 02.5, 025),
  c(09, 1.50, 1.50, 0.0, 1, 05, 014.00, 1.42857, 0.70000, 1, 1, 0.01, 02.5, 025),
  c(10, 1.50, 1.50, 0.0, 1, 05, 025.00, 0.80000, 1.25000, 1, 1, 0.01, 02.5, 025),
  c(11, 1.50, 1.50, 0.0, 1, 05, 021.00, 0.95238, 1.05000, 1, 1, 0.01, 02.5, 025),
  c(12, 1.50, 1.50, 0.0, 1, 05, 020.00, 1.00000, 1.00000, 1, 1, 0.01, 02.5, 025),
  c(13, 1.50, 1.50, 0.0, 1, 05, 020.00, 1.05000, 0.95238, 1, 1, 0.01, 02.5, 025),
  c(14, 1.50, 1.50, 0.0, 1, 05, 020.00, 1.25000, 0.80000, 1, 1, 0.01, 02.5, 025),
  c(15, 1.50, 1.50, 0.0, 1, 10, 020.00, 0.70000, 1.42857, 1, 1, 0.01, 02.5, 025),
  c(16, 1.50, 1.50, 0.0, 1, 10, 017.50, 0.80000, 1.25000, 1, 1, 0.01, 02.5, 025),
  c(17, 1.50, 1.50, 0.0, 1, 10, 015.56, 0.90000, 1.11111, 1, 1, 0.01, 02.5, 025),
  c(18, 1.50, 1.50, 0.0, 1, 10, 014.70, 0.95238, 1.05000, 1, 1, 0.01, 02.5, 025),
  c(19, 1.50, 1.50, 0.0, 1, 10, 014.00, 1.00000, 1.00000, 1, 1, 0.01, 02.5, 025),
  c(20, 1.50, 1.50, 0.0, 1, 10, 014.00, 1.05000, 0.95238, 1, 1, 0.01, 02.5, 025),
  c(21, 1.50, 1.50, 0.0, 1, 10, 014.00, 1.11111, 0.90000, 1, 1, 0.01, 02.5, 025),
  c(22, 1.50, 1.50, 0.0, 1, 10, 014.00, 1.25000, 0.80000, 1, 1, 0.01, 02.5, 025),
  c(23, 1.50, 1.50, 0.0, 1, 10, 014.00, 1.42857, 0.70000, 1, 1, 0.01, 02.5, 025),
  c(24, 1.50, 1.50, 0.0, 1, 10, 025.00, 0.80000, 1.25000, 1, 1, 0.01, 02.5, 025),
  c(25, 1.50, 1.50, 0.0, 1, 10, 021.00, 0.95238, 1.05000, 1, 1, 0.01, 02.5, 025),
  c(26, 1.50, 1.50, 0.0, 1, 10, 020.00, 1.00000, 1.00000, 1, 1, 0.01, 02.5, 025),
  c(27, 1.50, 1.50, 0.0, 1, 10, 020.00, 1.05000, 0.95238, 1, 1, 0.01, 02.5, 025),
  c(28, 1.50, 1.50, 0.0, 1, 10, 020.00, 1.25000, 0.80000, 1, 1, 0.01, 02.5, 025),
  c(29, 1.50, 1.50, 0.0, 0, 05, 014.70, 0.95238, 1.05000, 1, 1, 0.01, 02.5, 025),
  c(30, 1.50, 1.50, 0.0, 0, 05, 014.00, 1.05000, 0.95238, 1, 1, 0.01, 02.5, 025),
  c(31, 1.50, 1.50, 0.0, 0, 05, 021.00, 0.95238, 1.05000, 1, 1, 0.01, 02.5, 025),
  c(32, 1.50, 1.50, 0.0, 0, 05, 020.00, 1.05000, 0.95238, 1, 1, 0.01, 02.5, 025),
  c(33, 1.50, 1.50, 0.0, 0, 10, 014.70, 0.95238, 1.05000, 1, 1, 0.01, 02.5, 025),
  c(34, 1.50, 1.50, 0.0, 0, 10, 014.00, 1.05000, 0.95238, 1, 1, 0.01, 02.5, 025),
  c(35, 1.50, 1.50, 0.0, 0, 10, 021.00, 0.95238, 1.05000, 1, 1, 0.01, 02.5, 025),
  c(36, 1.50, 1.50, 0.0, 0, 10, 020.00, 1.05000, 0.95238, 1, 1, 0.01, 02.5, 025),
  c(37, 1.50, 1.50, 0.4, 1, 05, 020.00, 0.70000, 1.42857, 1, 2, 0.01, 02.5, 025),
  c(38, 1.50, 1.50, 0.4, 1, 05, 017.50, 0.80000, 1.25000, 1, 2, 0.01, 02.5, 025),
  c(39, 1.50, 1.50, 0.4, 1, 05, 015.56, 0.90000, 1.11111, 1, 2, 0.01, 02.5, 025),
  c(40, 1.50, 1.50, 0.4, 1, 05, 014.70, 0.95238, 1.05000, 1, 2, 0.01, 02.5, 025),
  c(41, 1.50, 1.50, 0.4, 1, 05, 014.00, 1.00000, 1.00000, 1, 2, 0.01, 02.5, 025),
  c(42, 1.50, 1.50, 0.4, 1, 05, 014.00, 1.05000, 0.95238, 1, 2, 0.01, 02.5, 025),
  c(43, 1.50, 1.50, 0.4, 1, 05, 014.00, 1.11111, 0.90000, 1, 2, 0.01, 02.5, 025),
  c(44, 1.50, 1.50, 0.4, 1, 05, 014.00, 1.25000, 0.80000, 1, 2, 0.01, 02.5, 025),
  c(45, 1.50, 1.50, 0.4, 1, 05, 014.00, 1.42857, 0.70000, 1, 2, 0.01, 02.5, 025),
  c(46, 1.50, 1.50, 0.4, 1, 05, 025.00, 0.80000, 1.25000, 1, 2, 0.01, 02.5, 025),
  c(47, 1.50, 1.50, 0.4, 1, 05, 021.00, 0.95238, 1.05000, 1, 2, 0.01, 02.5, 025),
  c(48, 1.50, 1.50, 0.4, 1, 05, 020.00, 1.00000, 1.00000, 1, 2, 0.01, 02.5, 025),
  c(49, 1.50, 1.50, 0.4, 1, 05, 020.00, 1.05000, 0.95238, 1, 2, 0.01, 02.5, 025),
  c(50, 1.50, 1.50, 0.4, 1, 05, 020.00, 1.25000, 0.80000, 1, 2, 0.01, 02.5, 025),
  c(51, 1.50, 1.50, 0.4, 1, 10, 020.00, 0.70000, 1.42857, 1, 2, 0.01, 02.5, 025),
  c(52, 1.50, 1.50, 0.4, 1, 10, 017.50, 0.80000, 1.25000, 1, 2, 0.01, 02.5, 025),
  c(53, 1.50, 1.50, 0.4, 1, 10, 015.56, 0.90000, 1.11111, 1, 2, 0.01, 02.5, 025),
  c(54, 1.50, 1.50, 0.4, 1, 10, 014.70, 0.95238, 1.05000, 1, 2, 0.01, 02.5, 025),
  c(55, 1.50, 1.50, 0.4, 1, 10, 014.00, 1.00000, 1.00000, 1, 2, 0.01, 02.5, 025),
  c(56, 1.50, 1.50, 0.4, 1, 10, 014.00, 1.05000, 0.95238, 1, 2, 0.01, 02.5, 025),
  c(57, 1.50, 1.50, 0.4, 1, 10, 014.00, 1.11111, 0.90000, 1, 2, 0.01, 02.5, 025),
  c(58, 1.50, 1.50, 0.4, 1, 10, 014.00, 1.25000, 0.80000, 1, 2, 0.01, 02.5, 025),
  c(59, 1.50, 1.50, 0.4, 1, 10, 014.00, 1.42857, 0.70000, 1, 2, 0.01, 02.5, 025),
  c(60, 1.50, 1.50, 0.4, 1, 10, 025.00, 0.80000, 1.25000, 1, 2, 0.01, 02.5, 025),
  c(61, 1.50, 1.50, 0.4, 1, 10, 021.00, 0.95238, 1.05000, 1, 2, 0.01, 02.5, 025),
  c(62, 1.50, 1.50, 0.4, 1, 10, 020.00, 1.00000, 1.00000, 1, 2, 0.01, 02.5, 025),
  c(63, 1.50, 1.50, 0.4, 1, 10, 020.00, 1.05000, 0.95238, 1, 2, 0.01, 02.5, 025),
  c(64, 1.50, 1.50, 0.4, 1, 10, 020.00, 1.25000, 0.80000, 1, 2, 0.01, 02.5, 025),
  c(65, 1.50, 1.50, 0.4, 0, 05, 014.70, 0.95238, 1.05000, 1, 2, 0.01, 02.5, 025),
  c(66, 1.50, 1.50, 0.4, 0, 05, 014.00, 1.05000, 0.95238, 1, 2, 0.01, 02.5, 025),
  c(67, 1.50, 1.50, 0.4, 0, 05, 021.00, 0.95238, 1.05000, 1, 2, 0.01, 02.5, 025),
  c(68, 1.50, 1.50, 0.4, 0, 05, 020.00, 1.05000, 0.95238, 1, 2, 0.01, 02.5, 025),
  c(69, 1.50, 1.50, 0.4, 0, 10, 014.70, 0.95238, 1.05000, 1, 2, 0.01, 02.5, 025),
  c(70, 1.50, 1.50, 0.4, 0, 10, 014.00, 1.05000, 0.95238, 1, 2, 0.01, 02.5, 025),
  c(71, 1.50, 1.50, 0.4, 0, 10, 021.00, 0.95238, 1.05000, 1, 2, 0.01, 02.5, 025),
  c(72, 1.50, 1.50, 0.4, 0, 10, 020.00, 1.05000, 0.95238, 1, 2, 0.01, 02.5, 025),
  c(73, 1.50, 1.50, 0.0, 1, 10, 150.00, 0.80000, 1.25000, 1, 3, 0.10, 10.0, 160),
  c(74, 1.50, 1.50, 0.0, 1, 10, 126.00, 0.95238, 1.05000, 1, 3, 0.10, 10.0, 160),
  c(75, 1.50, 1.50, 0.0, 1, 10, 120.00, 1.00000, 1.00000, 1, 3, 0.10, 10.0, 160),
  c(76, 1.50, 1.50, 0.0, 1, 10, 120.00, 1.05000, 0.95238, 1, 3, 0.10, 10.0, 160),
  c(77, 1.50, 1.50, 0.0, 1, 10, 120.00, 1.25000, 0.80000, 1, 3, 0.10, 10.0, 160),
  c(78, 1.50, 1.50, 0.4, 1, 10, 150.00, 0.80000, 1.25000, 1, 4, 0.10, 10.0, 160),
  c(79, 1.50, 1.50, 0.4, 1, 10, 126.00, 0.95238, 1.05000, 1, 4, 0.10, 10.0, 160),
  c(80, 1.50, 1.50, 0.4, 1, 10, 120.00, 1.00000, 1.00000, 1, 4, 0.10, 10.0, 160),
  c(81, 1.50, 1.50, 0.4, 1, 10, 120.00, 1.05000, 0.95238, 1, 4, 0.10, 10.0, 160),
  c(82, 1.50, 1.50, 0.4, 1, 10, 120.00, 1.25000, 0.80000, 1, 4, 0.10, 10.0, 160)
), nrow = 82, byrow = TRUE)
colnames(budget_restritions) <- c(
  "DecNum", "BasePayOne", "BasePayTwo", "Curvature", "FED", "Delay", "Budget",
  "PriceOneDivPriceTwo", "PriceTwoDivPriceOne", "ProbOne", "Block", "StepSize",
  "TickDist", "PlotMax"
)

# Generate data frame so that we can include all possible allocations that could be selected
# from the different budget restrictions
budget_restritions_df <- as.data.frame(budget_restritions)
# Add necessary columns, initialized with NA
budget_restritions_df$c_1_series <- list(NA)
budget_restritions_df$c_2_series <- list(NA)
budget_restritions_df$C_series <- list(NA)
budget_restritions_df$lb <- NA
budget_restritions_df$ub <- NA

# In the actual experiment, the sooner payment was displayed on the vertical axis, and the later
# payment was displayed on the horizontal axis. For convenience, we display c_1 on the horizontal
# and c_2 on the vertical axis here.

# Populate data frame
for (i in 1:dim(budget_restritions_df)[1]) {
  c_1_series <- seq(
    0, budget_restritions_df$Budget[i], budget_restritions_df$StepSize[i]
  )
  c_2_series <- round(c_2(
    c_1_series,
    budget_restritions_df$Budget[i],
    budget_restritions_df$PriceOneDivPriceTwo[i],
    budget_restritions_df$Curvature[i],
    budget_restritions_df$BasePayOne[i],
    budget_restritions_df$BasePayTwo[i]
  ), 2)  # Rounding to 2 decimal places because we displayed amounts with a precision of €0.01
  # Keep only c_1-c_2 pairs for c_2 >= baseline payment
  c_1_series <- c_1_series[c_2_series >= budget_restritions_df$BasePayTwo[i]]
  c_2_series <- c_2_series[1:length(c_1_series)]
  # Keep only c_1-c_2 pairs for c_1 >= baseline payment
  c_2_series <- c_2_series[c_1_series >= budget_restritions_df$BasePayOne[i]]
  c_1_series <- c_1_series[c_1_series >= budget_restritions_df$BasePayOne[i]]
  # Add horizontal/vertical segment at the position of the baseline payments
  c_1_series <- c(0, c_1_series, max(c_1_series))
  c_2_series <- c(max(c_2_series), c_2_series, 0)
  budget_restritions_df[[i, "c_1_series"]] <- as.list(c_1_series)
  budget_restritions_df[[i, "c_2_series"]] <- as.list(c_2_series)
  budget_restritions_df[[i, "C_series"]] <- as.list(c_1_series / c_2_series)
  budget_restritions_df[[i, "lb"]] <- budget_restritions_df[[i, "BasePayTwo"]] / max(c_2_series)
  budget_restritions_df[[i, "ub"]] <- max(c_1_series) / budget_restritions_df[[i, "BasePayTwo"]]
}
rm(c_1_series, c_2_series, i)




# PLOT BUDGET RESTRICTIONS --------------------------------------------------------------------


plotBRs <- function(df, curv) {
  df_aux <- df[df$Curvature == curv, ]
  plot_max <- max(df_aux[, "PlotMax"])
  plot(
    NaN, NaN,
    xlim = c(0, plot_max),
    ylim = c(0, plot_max),
    axes = FALSE,
    xlab = "", ylab = "",
    asp = 1,
    main = bquote("Budget restricions with curvature" ~ italic(z) ~ "=" ~ .(curv))
  )
  axis(1, seq(0, plot_max, 20), col = "gray25", col.axis = "gray25", pos = 0)
  text(x = plot_max / 2, y = -22.5, bquote(italic(c)[italic(t)]), xpd = TRUE)
  axis(2, seq(0, plot_max, 20), col = "gray25", col.axis = "gray25", pos = 0)
  text(x = -22.5, y = plot_max / 2, bquote(italic(c)[italic(t) + italic(k)]), xpd = TRUE, srt = 90)
  for (i in 1:dim(df_aux)[1]) {
    lines( unlist(df_aux[[i, "c_1_series"]]), unlist(df_aux[[i, "c_2_series"]]), col = "navy")
  }
}

for (curv in unique(budget_restritions_df$Curvature)) {
  plotBRs(budget_restritions_df, curv)
}
rm(curv)




# SIMULATE CHOICES ----------------------------------------------------------------------------


# For replicability of the results, use particular seed for the pseudorandom draws
set.seed(42)
# For new draws, remove the fixed seed via the following line
# set.seed(NULL)

# Optimal payment ratio:
# Equation (25) in https://www.econtribute.de/RePEc/ajk/ajkdps/ECONtribute_336_2024.pdf
C_star <- function(t, k, R, z, beta, delta, rho) {
  (1 / (beta^(t == 0) * delta^k * R^(1 + z)))^(1 / (rho + z))
}

# Add noise with Gaussian distribution:
# additive normally distributed noise on log(C_star) => multiplicative log-normal noise on C_star
C_star_noisy <- function(t, k, R, z, beta, delta, rho, sigma) {
  C_star(t, k, R, z, beta, delta, rho) * exp(rnorm(length(t), mean = 0, sd = sigma))
}

# Determine optimal (noisy) points on all BRs
opt_noisy_points_on_BR <- function(df_ind, beta, delta, rho, sigma) {
  ones <- rep(1, dim(df_ind)[1])
  beta_vec <- beta * ones
  delta_vec <- delta * ones
  rho_vec <- rho * ones
  sigma_vec <- sigma * ones
  # If rho is so small (negative) that the curvature of the utility function exceeds the curvature
  # of the budget restriction, replace rho by a value that is ever so slightly larger than
  # the curvature of the budget restriction. This way, the condition for an interior allocation
  # can still be applied, ideally leading to finite values (instead of NaNs), which can then be
  # converted to corner allocations.
  rho_vec[rho_vec <= -df_ind$Curvature] <-
    -df_ind$Curvature[rho_vec <= -df_ind$Curvature] + 0.000001
  C_star_noisy_list <- C_star_noisy(
    df_ind$FED,
    df_ind$Delay,
    df_ind$PriceOneDivPriceTwo,
    df_ind$Curvature,
    beta_vec, delta_vec, rho_vec, sigma_vec
  )
  # Initialize vectors with NA
  opt_noisy_C <- opt_noisy_C_idx <- opt_noisy_c_1 <- opt_noisy_c_2 <-
    rep(NA, length(C_star_noisy_list))
  # Populate vectors by iterating of the budget restrictions
  for (i in 1:length(C_star_noisy_list)) {
    # Find the point among the discrete points on the current budget restriction that is
    # closest to the (continuous) theoretical prediction
    dist <- abs(C_star_noisy_list[i] - unlist(df_ind[i, "C_series"]))
    # Use this allocation as the simulated choice
    opt_noisy_C_idx[i] <- which.min(dist)
    opt_noisy_C[i] <- unlist(df_ind[i, "C_series"])[opt_noisy_C_idx[i]]
    opt_noisy_c_1[i] <- unlist(df_ind[i, "c_1_series"])[opt_noisy_C_idx[i]]
    opt_noisy_c_2[i] <- unlist(df_ind[i, "c_2_series"])[opt_noisy_C_idx[i]]
    # Convert predicted allocation beyond the baseline payments to corner allocation
    max_C <- unlist(df_ind[i, "C_series"])[length(unlist(df_ind[i, "C_series"])) - 1]
    min_C <- unlist(df_ind[i, "C_series"])[2]
    if (opt_noisy_C[i] > max_C || is.nan(min(dist))) {
      opt_noisy_C[i] <- df_ind[[i, "ub"]]
      opt_noisy_c_1[i] <- max(unlist(df_ind[[i, "c_1_series"]]))
      opt_noisy_c_2[i] <- df_ind[[i, "BasePayTwo"]]
    }
    if (opt_noisy_C[i] < min_C) {
      opt_noisy_C[i] <- df_ind[[i, "lb"]]
      opt_noisy_c_1[i] <- df_ind[[i, "BasePayOne"]]
      opt_noisy_c_2[i] <- max(unlist(df_ind[[i, "c_2_series"]]))
    }
  }
  rm(dist)
  return(c(
    "C" = list(opt_noisy_C),
    "c_1" = list(opt_noisy_c_1),
    "c_2" = list(opt_noisy_c_2)
  ))
}

# Simulate participants with the following preference and noise parameters
# id, beta, delta, rho, sigma
params_sim <- matrix(c(
  1, 1.00, 1.00, 0.000, 0.00,  # Will be hard to estimate with z = 0 but easy with z = 0.4
  2, 1.00, 0.99, 0.000, 0.25,  # Will be hard to estimate with z = 0 but easy with z = 0.4
  3, 0.80, 0.99, 0.000, 0.25,  # Will be hard to estimate with z = 0 but easy with z = 0.4
  4, 0.80, 0.99, 0.100, 0.50,  # Should be estimable with both z = 0 and z = 0.4
  5, 0.85, 0.95, 0.150, 0.50,  # Should be estimable with both z = 0 and z = 0.4
  6, 1.00, 1.00, 1.000, 0.00,  # Should be estimable with both z = 0 and z = 0.4
  7, 1.00, 0.99, 100.0, 0.25,  # Should be estimable with both z = 0 and z = 0.4
  8, 1.00, 1.00, 200.0, 0.05,  # Should be estimable with both z = 0 and z = 0.4
  9, 1.00, 1.00, 200.0, 0.00   # May fail to converge (rho -> Inf), since c_1 = c_2 for all choices
), ncol = 5, byrow = TRUE)
colnames(params_sim) <- c("id", "beta", "delta", "rho", "sigma")

ids <- unique(params_sim[, "id"])

# Create new empty data frame
df <- budget_restritions_df[FALSE, ]

# Populate the data frame with the (noisy) choices of the simulated participants and
# plot the simulated choices
for (id in ids) {
  df_ind <- budget_restritions_df
  # Add column with IDs
  df_ind$id <- id
  # Store simulated choices
  points_sim <- opt_noisy_points_on_BR(
    df_ind,
    params_sim[params_sim[, "id"] == id, "beta"],
    params_sim[params_sim[, "id"] == id, "delta"],
    params_sim[params_sim[, "id"] == id, "rho"],
    params_sim[params_sim[, "id"] == id, "sigma"]
  )
  df_ind$payment_1 <- unlist(points_sim["c_1"])
  df_ind$payment_2 <- unlist(points_sim["c_2"])
  df_ind$payment_ratio <- unlist(points_sim["C"])
  # Append to the existing data frame
  df <- rbind(df, df_ind)
  # Generate a plot per level of curvature of the budget restrictions
  for (curv in unique(df_ind$Curvature)) {
    plotBRs(df_ind, curv)
    points(
      x = unlist(points_sim["c_1"])[df_ind$Curvature == curv],
      y = unlist(points_sim["c_2"])[df_ind$Curvature == curv],
      col = "navy",
      pch = 20,
    )
    mtext(bquote(
      "Simulated choices for ID" ~ .(id) * ":" ~
        italic(β) ~ "=" ~ .(params_sim[params_sim[, "id"] == id, "beta"]) * "," ~
        italic(δ) ~ "=" ~ .(params_sim[params_sim[, "id"] == id, "delta"]) * "," ~
        italic(ρ) ~ "=" ~ .(params_sim[params_sim[, "id"] == id, "rho"]) * "," ~
        italic(σ) ~ "=" ~ .(params_sim[params_sim[, "id"] == id, "sigma"])
    ), side = 3, col = "navy", line = 0.25)
    Sys.sleep(0.5)
  }
}
rm(curv, df_ind, id)
rm(points_sim)




# TOBIT ESTIMATION ----------------------------------------------------------------------------


# Generating the explanatory variables, see
# https://www.econtribute.de/RePEc/ajk/ajkdps/ECONtribute_336_2024.pdf, eq. (26) on p. 15:
df$cov_1 <-
  -as.integer(df$FED == 0)  # -I[t = 0], coefficient: gamma_beta
df$cov_2 <-
  -df$Delay  # -k, coefficient: gamma_delta
df$cov_3 <-
  -(1 + df$Curvature) * log(df$PriceOneDivPriceTwo)  # -(1 + z) ln(R), coefficient: gamma_rho

ids <- unique(df$id)
tobit_estimates_report <- list()

# Estimate separately for each curvature level of the budget restrictions
for (curv in unique(df$Curvature)) {
  # Initialize collection of estimates with NA
  tobit_estimates <- matrix(rep(NA, length(ids) * 5), ncol = 5)
  tobit_estimates <- cbind(ids, tobit_estimates)
  colnames(tobit_estimates) <- c("id", "beta", "delta", "rho", "sigma", "logL")
  # Estimate each individual separately
  for (id in ids) {
    print(paste0("Tobit estimation, Subject ID: ", id, "; BR curvature: ", curv))
    df_ind <- df[df$id == id & df$Curvature == curv, ]
    # Plot simulated choices
    plotBRs(df_ind, curv)
    points(
      df_ind$payment_1,
      df_ind$payment_2,
      col = "navy",
      pch = 20
    )
    mtext(bquote(
      "Simulated choices for ID" ~ .(id) * ":" ~
        italic(β) ~ "=" ~ .(params_sim[params_sim[, "id"] == id, "beta"]) * "," ~
        italic(δ) ~ "=" ~ .(params_sim[params_sim[, "id"] == id, "delta"]) * "," ~
        italic(ρ) ~ "=" ~ .(params_sim[params_sim[, "id"] == id, "rho"]) * "," ~
        italic(σ) ~ "=" ~ .(params_sim[params_sim[, "id"] == id, "sigma"])
    ), side = 3, col = "navy", line = 0.25)
    tryCatch(
      {
        withCallingHandlers(
          # Attempt Tobit regression
          {
            model <- tobit(
              log(payment_ratio) ~ cov_1 + cov_2 + cov_3 - 1,
              # "+" here because we included the minus sign above when creating the regressors
              left = log(df_ind$lb),
              right = log(df_ind$ub),
              data = df_ind
            )
          },
          # Show potential warning messages of the Tobit regression
          warning = function(w) {
            message(style_bold(col_blue("WARNING: ", conditionMessage(w))))
            invokeRestart("muffleWarning")
          }
        )
        vcov_matrix <- vcov(model)
        coeffs <- model$coefficients
        est_beta <-
          round(exp(as.numeric(coeffs[["cov_1"]]) / as.numeric(coeffs[["cov_3"]])), round_prec)
          # beta = exp(gamma_beta / gamma_rho), eq. (27) in Gerhardt & Suchy (2024)
        est_delta <-
          round(exp(as.numeric(coeffs[["cov_2"]]) / as.numeric(coeffs[["cov_3"]])), round_prec)
          # delta = exp(gamma_delta / gamma_rho), eq. (28)
        est_rho <-
          round(1 / as.numeric(coeffs[["cov_3"]]) - curv, round_prec)
          # rho = (1 / gamma_rho) - z, eq. (29)
        est_sigma <- round(model$scale, round_prec)
        tobit_estimates[tobit_estimates[, "id"] == id, "beta"] <- est_beta
        tobit_estimates[tobit_estimates[, "id"] == id, "delta"] <- est_delta
        tobit_estimates[tobit_estimates[, "id"] == id, "rho"] <- est_rho
        tobit_estimates[tobit_estimates[, "id"] == id, "sigma"] <- est_sigma
        tobit_estimates[tobit_estimates[, "id"] == id, "logL"] <- round(logLik(model), round_prec)
        # Use delta method to calculate standard errors of structural parameters
        est_beta_se <- deltamethod(list(~ exp(x1 / x3)), coeffs, vcov(model)[1:3, 1:3])
        est_delta_se <- deltamethod(list(~ exp(x2 / x3)), coeffs, vcov(model)[1:3, 1:3])
        est_rho_se <- deltamethod(list(~ 1 / x3 - curv), coeffs, vcov(model)[1:3, 1:3])
        # Add best-fitting allocations to plot (remove random component, i.e., sigma = 0)
        points(
          unlist(opt_noisy_points_on_BR(df_ind, est_beta, est_delta, est_rho, sigma = 0)["c_1"]),
          unlist(opt_noisy_points_on_BR(df_ind, est_beta, est_delta, est_rho, sigma = 0)["c_2"]),
          col = "#FFA50099",
          pch = 20
        )
        mtext(bquote(
          "Parameter estimates:" ~
            hat(italic(β)) ~ "=" ~ .(est_beta) * "," ~
            hat(italic(δ)) ~ "=" ~ .(est_delta) * "," ~
            hat(italic(ρ)) ~ "=" ~ .(est_rho) * "," ~
            hat(italic(σ)) ~ "=" ~ .(est_sigma)
        ), side = 3, col = "#FFA500", line = -1)
      },
      # If Tobit regression fails, issue error message
      error = function(e) {
        message(style_bold(bg_red(col_br_white("ERROR: ", conditionMessage(e)))))
      }
    )
    Sys.sleep(0.25)  # Short break to update plot
  }
  # Collect and display Tobit estimates
  tobit_estimates_report[[toString(curv)]] <- tobit_estimates
  print(cbind(params_sim, tobit_estimates_report[[toString(curv)]][, 2:6]))
}
# Collect estimates
all_estimates_report <- tobit_estimates_report
rm(curv, id, df_ind, tobit_estimates)
rm(coeffs, vcov_matrix)
rm(est_beta, est_beta_se, est_delta, est_delta_se, est_rho, est_rho_se, est_sigma)




# NONLINEAR MAXIMUM LIKELIHOOD ESTIMATION -----------------------------------------------------


# Log-likelihood contribution of a single observation according to eq. (30) in
# https://www.econtribute.de/RePEc/ajk/ajkdps/ECONtribute_336_2024.pdf
LL_contrib <- function(C_star_obs, t, k, lb, ub, R, z, beta, delta, rho, sigma) {
  C_star_pred <- C_star(t, k, R, z, beta, delta, rho)
  if (
    beta < 0 || delta < 0 || sigma < 0 || rho < -z
    # Parameters must not become negative, and
    # curvature of utility beyond curvature of BR cannot be identified
  ) {
    lnf <- -999999
  } else {
    # Interior solution
    lnf <- log(dnorm((log(C_star_obs) - log(C_star_pred)) / sigma) / sigma)
    # If observation is lower bound
    if (C_star_obs < lb + 0.00001) {
      lnf <- log(pnorm((log(lb) - log(C_star_pred)) / sigma))
    }
    # If observation is upper bound
    if (C_star_obs > ub - 0.00001) {
      lnf <- log(pnorm((log(C_star_pred) - log(ub)) / sigma))
    }
  }
  # Rule out NaNs and infinite values
  if (is.na(lnf) || is.nan(lnf) || lnf == -Inf) {
    lnf <- -999999
  }
  return(lnf)
}

# Log-likelihood contributions of all observations collected in vector
# (required by some optimization methods)
LL_contrib_vec <- function(C_star_obs, t, k, lb, ub, R, z, beta, delta, rho, sigma) {
  ln <- unlist(sapply(
    1:length(C_star_obs),
    function(i) {
      LL_contrib(
        C_star_obs[i],
        t[i], k[i], lb[i], ub[i], R[i], z[i],
        beta, delta, rho, sigma
      )
    }
  ))
  return(as.vector(ln))
}

ids <- unique(df$id)
mle_estimates_report <- list()

# Estimate separately for each curvature level of the budget restrictions
for (curv in unique(df$Curvature)) {
  # Initialize collection of estimates with NA
  mle_estimates <- matrix(rep(NA, length(ids) * 5), ncol = 5)
  mle_estimates <- cbind(ids, mle_estimates)
  colnames(mle_estimates) <- c("id", "beta", "delta", "rho", "sigma", "logL")
  # Estimate each individual separately
  for (id in ids) {
    print(paste0("NL-ML estimation, Subject ID: ", id, "; BR curvature: ", curv))
    df_ind <- df[df$id == id & df$Curvature == curv, ]
    # Objective function for maxLik must not contain any arguments except params
    LL_contrib_vec_filled <- function(params) {
      LL_contrib_vec(
        df_ind$payment_ratio,
        df_ind$FED, df_ind$Delay, df_ind$lb, df_ind$ub,
        df_ind$PriceOneDivPriceTwo, df_ind$Curvature,
        params[1], params[2], params[3], params[4]
      )
    }
    # Set the initial values for the numerical nonlinear estimation procedure
    # By default, take the outcome of the Tobit regression
    tobit_for_init <- tobit_estimates_report[[toString(curv)]]
    init_vals <- round(c(
      tobit_for_init[tobit_for_init[, "id"] == id, "beta"],
      tobit_for_init[tobit_for_init[, "id"] == id, "delta"],
      tobit_for_init[tobit_for_init[, "id"] == id, "rho"],
      tobit_for_init[tobit_for_init[, "id"] == id, "sigma"]
    ), 2)
    # If the Tobit regression yields nonsensical results, set different initical value
    if (init_vals["beta"] <= 0 | is.na(init_vals["beta"])) {
      init_vals["beta"] <- 0.95
    }
    if (init_vals["delta"] <= 0 | is.na(init_vals["delta"])) {
      init_vals["delta"] <- 0.95
    }
    if (is.na(init_vals["rho"])) {
      init_vals["rho"] <- 0.05
    } else if (init_vals["rho"] <= -curv) {
      init_vals["rho"] <- 1
    }
    if (init_vals["sigma"] <= 0 | is.na(init_vals["sigma"])) {
      init_vals["sigma"] <- 0.1
    }
    # As is frequently the case with numerical optimization procedures, individual-specific
    # initial values may be required, e.g., if Tobit did not converge. This is particularly likely
    # with linear budget restrictions and less so with strictly concave budget restrictions.
    if (id == 1 && curv == 0) {
      init_vals[c("beta", "delta", "rho", "sigma")] = c(1, 1, 0.01, 0.005)
    }
    if (id %in% c(2, 3) && curv == 0) {
      init_vals[c("beta", "delta", "rho", "sigma")] = c(0.995, 0.995, 0.01, 0.2)
    }
    # sum(LL_contrib_vec_filled(init_vals))  # Helpful for finding initial values
    # Plot simulated choices
    plotBRs(df_ind, curv)
    points(
      df_ind$payment_1,
      df_ind$payment_2,
      col = "navy",
      pch = 20
    )
    mtext(bquote(
      "Simulated choices for ID" ~ .(id) * ":" ~
        italic(β) ~ "=" ~ .(params_sim[params_sim[, "id"] == id, "beta"]) * "," ~
        italic(δ) ~ "=" ~ .(params_sim[params_sim[, "id"] == id, "delta"]) * "," ~
        italic(ρ) ~ "=" ~ .(params_sim[params_sim[, "id"] == id, "rho"]) * "," ~
        italic(σ) ~ "=" ~ .(params_sim[params_sim[, "id"] == id, "sigma"])
    ), side = 3, col = "navy", line = 0.25)
    tryCatch(
      {
        withCallingHandlers(
          # Attempt NL-MLE
          {
            mle_estim <- maxLik(
              logLik = LL_contrib_vec_filled,
              start = init_vals,
              # method = "BFGS",  # Broyden/Fletcher/Goldfarb/Shanno
              # method = "BFGSR",  # Broyden/Fletcher/Goldfarb/Shanno
              # method = "BHHH",  # Berndt/Hall/Hall/Hausman
              method = "NR", # Newton/Raphson
              # method = "SANN",  # Simulated Annealing
              control = list(
                gradtol = 10^-8, # Return code 1 (normal convergence)
                tol = 10^-8, # Return code 2 (normal convergence)
                steptol = -1, # Return code 3
                iterlim = 1000, # 10^6,  # Return code 4
                reltol = 10^-8 # Return code 8 (normal convergence)
              )
            )
            est_se <- stdEr(mle_estim, eigentol = 10^(-15))
          },
          # Show potential warning messages of the NL-MLE
          warning = function(w) {
            message(style_bold(col_blue("WARNING: ", conditionMessage(w))))
            invokeRestart("muffleWarning")
          }
        )
        # Add best-fitting allocations to plot (remove random component, i.e., sigma = 0)
        points(
          unlist(opt_noisy_points_on_BR(
            df_ind,
            mle_estim$estimate["beta"], mle_estim$estimate["delta"], mle_estim$estimate["rho"],
            sigma = 0
          )["c_1"]),
          unlist(opt_noisy_points_on_BR(
            df_ind,
            mle_estim$estimate["beta"], mle_estim$estimate["delta"], mle_estim$estimate["rho"],
            sigma = 0
          )["c_2"]),
          col = "#FF450088",
          pch = 20
        )
        mle_estim$estimate <- round(mle_estim$estimate, round_prec)
        mtext(bquote(
          "Parameter estimates:" ~
            hat(italic(β)) ~ "=" ~ .(mle_estim$estimate["beta"]) * "," ~
            hat(italic(δ)) ~ "=" ~ .(mle_estim$estimate["delta"]) * "," ~
            hat(italic(ρ)) ~ "=" ~ .(mle_estim$estimate["rho"]) * "," ~
            hat(italic(σ)) ~ "=" ~ .(mle_estim$estimate["sigma"])
        ), side = 3, col = "#FF4500", line = -1)
        mle_estimates[mle_estimates[, "id"] == id, "beta"] <- mle_estim$estimate["beta"]
        mle_estimates[mle_estimates[, "id"] == id, "delta"] <- mle_estim$estimate["delta"]
        mle_estimates[mle_estimates[, "id"] == id, "rho"] <- mle_estim$estimate["rho"]
        mle_estimates[mle_estimates[, "id"] == id, "sigma"] <- mle_estim$estimate["sigma"]
        mle_estimates[mle_estimates[, "id"] == id, "logL"] <- round(mle_estim$maximum, round_prec)
      },
      # If NL-MLE fails, issue error message
      error = function(e) {
        message(style_bold(bg_red(col_br_white("ERROR: ", conditionMessage(e)))))
      }
    )
    Sys.sleep(0.25)  # Short break to update plot
  }
  # Collect NL-MLE estimates
  mle_estimates_report[[toString(curv)]] <- mle_estimates
  cols = c("beta", "delta", "rho", "sigma", "logL")
  # Collect and display all estimates
  all_estimates_report[[toString(curv)]] <- cbind(
    params_sim,
    tobit_estimates_report[[toString(curv)]][, cols],
    mle_estimates_report[[toString(curv)]][, cols]
  )
  colnames(all_estimates_report[[toString(curv)]]) <-
    c(
      "id",
      "beta_sim", "delta_sim", "rho_sim", "sigma_sim",
      "beta_tobit", "delta_tobit", "rho_tobit", "sigma_tobit", "logL_tobit",
      "beta_mle", "delta_mle", "rho_mle", "sigma_mle", "logL_mle"
    )
  print(all_estimates_report[[toString(curv)]])
}
rm(curv, id, df_ind, mle_estimates, est_se, tobit_for_init)

all_estimates_report[["0"]][, 1:10]
# This should yield
#      id beta_sim delta_sim rho_sim sigma_sim beta_tobit delta_tobit rho_tobit sigma_tobit logL_tobit
# [1,]  1     1.00      1.00    0.00      0.00         NA          NA        NA      0.0006    32.1004
# [2,]  2     1.00      0.99    0.00      0.25     0.9557      0.9910    0.0000    155.4723 -9471.4181
# [3,]  3     0.80      0.99    0.00      0.25     0.9557      0.9910    0.0000    155.4723 -9471.4181
# [4,]  4     0.80      0.99    0.10      0.50     0.7998      0.9907    0.0961      0.3942   -15.9651
# [5,]  5     0.85      0.95    0.15      0.50     0.8554      0.9514    0.1556      0.4335   -14.5728
# [6,]  6     1.00      1.00    1.00      0.00     1.0000      1.0000    0.9999      0.0001   324.9009
# [7,]  7     1.00      0.99  100.00      0.25     0.6308      0.9941   -3.9738      0.1959     8.6606
# [8,]  8     1.00      1.00  200.00      0.05     0.1828      1.0957   57.7050      0.0496    64.9979
# [9,]  9     1.00      1.00  200.00      0.00     0.9814      1.0025  196.4053      0.0007   240.3492
# With linear budget restrictions:
# IDs 1, 2, and 3: Tobit does not converge for linear utility.
# Tobit converges to strongly convex utility instead of strongly concave utility for ID 7.
# ID 8: Strongly concave utility is hard to estimate in the presence of noise.

all_estimates_report[["0"]][, c(1:5, 11:15)]
# This should yield
#      id beta_sim delta_sim rho_sim sigma_sim beta_mle delta_mle  rho_mle sigma_mle logL_mle
# [1,]  1     1.00      1.00    0.00      0.00   1.0000    1.0000   0.0110    0.0003  34.8556
# [2,]  2     1.00      0.99    0.00      0.25   0.9905    0.9900   0.0001    0.1853   0.0000
# [3,]  3     0.80      0.99    0.00      0.25   0.9905    0.9900   0.0001    0.1853   0.0000
# [4,]  4     0.80      0.99    0.10      0.50   0.7998    0.9907   0.0961    0.3942 -15.9651
# [5,]  5     0.85      0.95    0.15      0.50   0.8554    0.9514   0.1556    0.4335 -14.5728
# [6,]  6     1.00      1.00    1.00      0.00   1.0000    1.0000   0.9999    0.0001 324.9009
# [7,]  7     1.00      0.99  100.00      0.25 555.2478    1.0991  57.9643    0.2008   7.6510
# [8,]  8     1.00      1.00  200.00      0.05   0.1745    1.0984  59.2789    0.0496  64.9978
# [9,]  9     1.00      1.00  200.00      0.00   0.9800    1.0000 196.4100    0.0049 179.5666
# With linear budget restrictions (LBRs):
# IDs 1, 2, and 3: Estimation becomes possible with NL-MLE by searching for suitable initial values,
# which, however, is cumbersome and difficult to automate.
# ID 7: Our NL-MLE routine converges to strongly concave utility, thereby fixing the weak point
# of Tobit. At the same time, with strongly concave utility all allocations are characterized by
# c_1 = c_2, with beta and delta hardly having any influence. Hence, identification of beta and
# delta becomes very hard.

all_estimates_report[["0.4"]][, 1:10]
# This should yield
#      id beta_sim delta_sim rho_sim sigma_sim beta_tobit delta_tobit rho_tobit sigma_tobit logL_tobit
# [1,]  1     1.00      1.00    0.00      0.00     1.0000      1.0000   -0.0002      0.0004   260.7333
# [2,]  2     1.00      0.99    0.00      0.25     0.9697      0.9903   -0.0093      0.2442    -0.3814
# [3,]  3     0.80      0.99    0.00      0.25     0.8085      0.9913    0.0063      0.2156     4.7404
# [4,]  4     0.80      0.99    0.10      0.50     0.8394      0.9889    0.0885      0.5054   -30.1979
# [5,]  5     0.85      0.95    0.15      0.50     0.9753      0.9509    0.0958      0.4442   -24.9066
# [6,]  6     1.00      1.00    1.00      0.00     0.9997      1.0000    0.9981      0.0004   258.3382
# [7,]  7     1.00      0.99  100.00      0.25     0.8767      0.9071   12.5038      0.2460    -0.6843
# [8,]  8     1.00      1.00  200.00      0.05     1.2020      1.0142   18.5775      0.0513    63.6310
# [9,]  9     1.00      1.00  200.00      0.00     0.9419      1.0080  196.3481      0.0006   248.9800
# With strictly concave budget restrictions (SCBRs):
# IDs 1, 2, and 3: Tobit converges for linear utility. Thus, SCBRs fix the weak point of LBRs.

all_estimates_report[["0.4"]][, c(1:5, 11:15)]
# This should yield
#      id beta_sim delta_sim rho_sim sigma_sim beta_mle delta_mle  rho_mle sigma_mle logL_mle
# [1,]  1     1.00      1.00    0.00      0.00   1.0000    1.0000  -0.0002    0.0004 260.7333
# [2,]  2     1.00      0.99    0.00      0.25   0.9697    0.9903  -0.0093    0.2442  -0.3814
# [3,]  3     0.80      0.99    0.00      0.25   0.8085    0.9913   0.0063    0.2156   4.7404
# [4,]  4     0.80      0.99    0.10      0.50   0.8394    0.9889   0.0885    0.5054 -30.1979
# [5,]  5     0.85      0.95    0.15      0.50   0.9753    0.9509   0.0958    0.4442 -24.9066
# [6,]  6     1.00      1.00    1.00      0.00   0.9997    1.0000   0.9981    0.0004 258.3382
# [7,]  7     1.00      0.99  100.00      0.25   0.8767    0.9071  12.5036    0.2460  -0.6843
# [8,]  8     1.00      1.00  200.00      0.05   1.2019    1.0142  18.5773    0.0513  63.6310
# [9,]  9     1.00      1.00  200.00      0.00   0.9400    1.0100 196.3500    0.0047 181.4055
# This illustrates that with the chosen optimization method ("NR"), our NL-MLE routine delivers
# (virtually) the same results as Tobit.
\end{lstlisting}

\clearpage

\iftrue
	% Use \iffalse to suppress the text font samples, which may be needed for compilation via https://uni-bonn.sciebo.de/apps/overleaf_nextcloud/launcher/launch
	% Use \iftrue to display the text font samples

	\begingroup

	\normalsize

	\section{Font Samples}
	\label{app:font-samples}

	\subsection[Font Sample Times]{\fontfamily{ntxtlf}Font Sample Times}
	\label{app:font-samples:times}

	\begin{otherlanguage}{ngerman}\noindent\normalsize%
		\fontfamily{ntxtlf}%
		\makeatletter%
		\ifdim \f@size pt < 11.5pt%
			\fontsize{11pt}{\baselinedist}%  % 10.945pt
		\else%
			\fontsize{12pt}{\baselinedist}%
		\fi%
		\makeatother%
		\selectfont%
		\Kafka
		\medskip
		\noindent\url{https://en.wikipedia.org/wiki/Times_New_Roman} \par
	\end{otherlanguage}

	\subsection[Font Sample Palatino]{\fontfamily{zpltlf}\selectfont Font Sample Palatino}
	\label{app:font-samples:palatino}

	\begin{otherlanguage}{ngerman}\noindent%
		\fontfamily{zpltlf}%
		\makeatletter%
		\ifdim \f@size pt < 11.5pt%
			\fontsize{10pt}{\baselinedist}%
		\else%
			\fontsize{10.95pt}{\baselinedist}%
		\fi%
		\makeatother%
		\selectfont%
		\Kafka
		\medskip
		\noindent\url{https://en.wikipedia.org/wiki/Palatino} \par
	\end{otherlanguage}

	\clearpage

	\section*{\strut}

	\subsection[Font Sample Utopia]{\fontfamily{erewhon-TLF}\selectfont Font Sample Utopia}
	\label{app:font-samples:utopia}

	\begin{otherlanguage}{ngerman}\noindent%
		\fontfamily{erewhon-TLF}%
		\makeatletter%
		\ifdim \f@size pt < 11.5pt%
			\fontsize{10.5pt}{\baselinedist}%
		\else%
			\fontsize{11.43pt}{\baselinedist}%
		\fi%
		\makeatother%
		\selectfont%
		\Kafka
		\medskip
		\noindent\url{https://en.wikipedia.org/wiki/Utopia_(typeface)} \par
	\end{otherlanguage}

	\subsection[Font Sample Charter]{\fontfamily{XCharter-TLF}\selectfont Font Sample Charter}
	\label{app:font-samples:charter}

	\begin{otherlanguage}{ngerman}\noindent%
		\fontfamily{XCharter-TLF}%
		\makeatletter%
		\ifdim \f@size pt < 11.5pt%
			\fontsize{10.26pt}{\baselinedist}%
		\else%
			\fontsize{11.1pt}{\baselinedist}%
		\fi%
		\makeatother%
		\selectfont%
		\Kafka
		\medskip
		\noindent\url{https://en.wikipedia.org/wiki/Bitstream_Charter} \par
	\end{otherlanguage}

	\clearpage

	\section*{\strut}

	\subsection[Font Sample {\textls[40]{STIX}} Two]{\fontfamily{SticksTooText-TLF}\selectfont Font Sample \textls[40]{STIX} Two}
	\label{app:font-samples:stix-two}

	\begin{otherlanguage}{ngerman}\noindent%
		\fontfamily{SticksTooText-TLF}%
		\makeatletter%
		\ifdim \f@size pt < 11.5pt%
			\fontsize{10.55pt}{\baselinedist}%
		\else%
			\fontsize{11.51pt}{\baselinedist}%
		\fi%
		\makeatother%
		\selectfont%
		\Kafka
		\medskip
		\noindent\url{https://en.wikipedia.org/wiki/STIX_Fonts_project#STIX_2.0.0} \par
	\end{otherlanguage}

	\subsection[Font Sample Libertinus Serif]{\fontfamily{LibertinusSerif-TLF}\selectfont Font Sample Libertinus Serif}
	\label{app:font-samples:libertinus-serif}

	\begin{otherlanguage}{ngerman}\noindent%
		\fontfamily{LibertinusSerif-TLF}%
		\makeatletter%
		\ifdim \f@size pt < 11.5pt%
			\fontsize{10.92pt}{\baselinedist}%
		\else%
			\fontsize{11.85pt}{\baselinedist}%
		\fi%
		\makeatother%
		\selectfont%
		\Kafka
		\medskip
		\noindent\url{https://en.wikipedia.org/wiki/Libertinus} \par
	\end{otherlanguage}

	\clearpage

	\section*{\strut}

	\subsection[Font Sample New Century Schoolbook]{\fontfamily{TeXGyreScholaX-TLF}\selectfont{\relscale{0.925}Font Sample New Century Schoolbook}}
	\label{app:font-samples:new-century-schoolbook}

	\begin{otherlanguage}{ngerman}\noindent%
		\fontfamily{TeXGyreScholaX-TLF}%
		\makeatletter%
		\ifdim \f@size pt < 11.5pt%
			\fontsize{9.6pt}{\baselinedist}%
		\else%
			\fontsize{10.57pt}{\baselinedist}%
		\fi%
		\makeatother%
		\selectfont%
		\Kafka
		\medskip
		\noindent\url{https://en.wikipedia.org/wiki/Century_type_family} \par
	\end{otherlanguage}

	\endgroup

	\clearpage

\fi

\iffalse
	% Use \iffalse to suppress the math font samples, needed for https://uni-bonn.sciebo.de/apps/overleaf_nextcloud/launcher/launch
	% Use \iftrue to display the math font samples

	\section{Math \texorpdfstring{``}{“}Torture Test\texorpdfstring{''}{”}}
	\label{app:math-torture}

	\makeatletter
	\newcommand*{\checkgreekletters}{%
		\@for\@tempa:=%
		alpha,beta,gamma,delta,epsilon,varepsilon,zeta,eta,theta,vartheta,iota,kappa,lambda,mu,nu,xi,%
		omicron,pi,varpi,rho,varrho,sigma,varsigma,tau,upsilon,phi,varphi,chi,psi,omega,digamma,%
		Alpha,Beta,Gamma,Delta,Epsilon,Zeta,Eta,Theta,Iota,Kappa,Lambda,Mu,Nu,Xi,%
		Omicron,Pi,Rho,Sigma,Tau,Upsilon,Phi,Chi,Psi,Omega,Digamma%
		\do{$\csname\@tempa\endcsname,$ }%
	}
	\newcommand*{\checkgreeklettersup}{%
		\@for\@tempa:=%
		upalpha,upbeta,upgamma,updelta,upepsilon,upvarepsilon,upzeta,upeta,uptheta,upvartheta,upiota,upkappa,uplambda,upmu,upnu,upxi,%
		upomicron,uppi,upvarpi,uprho,upvarrho,upsigma,upvarsigma,uptau,upupsilon,upphi,upvarphi,upchi,uppsi,upomega,updigamma,%
		upAlpha,upBeta,upGamma,upDelta,upEpsilon,upZeta,upEta,upTheta,upIota,upKappa,upLambda,upMu,upNu,upXi,%
		upOmicron,upPi,upRho,upSigma,upTau,upUpsilon,upPhi,upChi,upPsi,upOmega,upDigamma%
		\do{$\csname\@tempa\endcsname,$ }%
	}
	\makeatother

	\def\test#1{#1}

	\def\testnums{%
	  \test 0 \test 1 \test 2 \test 3 \test 4 \test 5 \test 6 \test 7
	  \test 8 \test 9 }
	\def\testupperi{%
	  \test A \test B \test C \test D \test E \test F \test G \test H
	  \test I \test J \test K \test L \test M }
	\def\testupperii{%
	  \test N \test O \test P \test Q \test R \test S \test T \test U
	  \test V \test W \test X \test Y \test Z }
	\def\testupper{%
	  \testupperi\testupperii}

	\def\testloweri{%
	  \test a \test b \test c \test d \test e \test f \test g \test h
	  \test i \test j \test k \test l \test m }
	\def\testlowerii{%
	  \test n \test o \test p \test q \test r \test s \test t \test u
	  \test v \test w \test x \test y \test z }
	\def\testlower{%
	  \testloweri\testlowerii}

	\def\testupgreeki{%
	  \test A \test B \test\Gamma \test\Delta \test E \test Z \test H
	  \test\Theta \test I \test K \test\Lambda \test M }
	\def\testupgreekii{%
	  \test N \test\Xi \test O \test\Pi \test P \test\Sigma \test T
	  \test\Upsilon \test\Phi \test X \test\Psi \test\Omega }
	\def\testupgreek{%
	  \testupgreeki\testupgreekii}

	\def\testlowgreeki{%
	  \test\alpha \test\beta \test\gamma \test\delta \test\epsilon
	  \test\zeta \test\eta \test\theta \test\iota \test\kappa \test\lambda
	  \test\mu }
	\def\testlowgreekii{%
	  \test\nu \test\xi \test o \test\pi \test\rho \test\sigma \test\tau
	  \test\upsilon \test\phi \test\chi \test\psi \test\omega }
	\def\testlowgreekiii{%
	  \test\varepsilon \test\vartheta \test\varpi \test\varrho
	  \test\varsigma \test\varphi}
	\def\testlowgreek{%
	  \testlowgreeki\testlowgreekii\testlowgreekiii}

	\def\testlowgreekupi{%
	  \test\upalpha \test\upbeta \test\upgamma \test\updelta \test\upepsilon
	  \test\upzeta \test\upeta \test\uptheta \test\upiota \test\upkappa \test\uplambda
	  \test\upmu }
	\def\testlowgreekupii{%
	  \test\upnu \test\upxi \test o \test\uppi \test\uprho \test\upsigma \test\uptau
	  \test\upupsilon \test\upphi \test\upchi \test\uppsi \test\upomega }
	\def\testlowgreekupiii{%
	  \test\upvarepsilon \test\upvartheta \test\upvarpi \test\upvarrho
	  \test\upvarsigma \test\upvarphi}
	\def\testlowgreekup{%
	  \testlowgreekupi\testlowgreekupii\testlowgreekupiii}

	\newcommand{\showfamily}{}

	Most of the following examples are taken from \textit{The \TeX book} \citep[][see \url{https://ctan.org/pkg/texbook}]{knuth:ct:a} and were adapted for \LaTeX\ from Karl Berry's torture test for plain \TeX\ math fonts.

	\noindent $x + y - z$, \quad $x + y * z$, \quad $z * y \divslash z$, \quad
	$(x+y)\,(x-y) = x^2 - y^2$,

	\noindent $x \times y \cdot z = [x\, y\, z]$, \quad $x\circ y \bullet z$, \quad
	$x\cup y \cap z$, \quad $x\sqcup y \sqcap z$, \quad

	\noindent $x \vee y \wedge z$, \quad $x\pm y\mp z$, \quad
	$x=y \divslash z$, \;\; $x \coloneqq y$, \;\; $x\le y \ne z$, \;\; $x \sim y \simeq z$
	$x \equiv y \nequiv z$, \;\; $x\subset y \subseteq z$

	\noindent $\sin2\theta=2\sin\theta\cos\theta$, \quad
	$\hbox{O}(n\log n\log n)$, \quad
	$\Pr(X>x)=\exp(-x \divslash \mu)$,

	\noindent $\bigl(x\in A(n)\bigm|x\in B(n)\bigr)$, \quad
	$\bigcup_n X_n\bigm\|\bigcap_n Y_n$

	% page 178

	\noindent In-text matrices $\binom{1\ 1}{0\ 1}$ and $\bigl(\genfrac{}{}{0pt}{}{a}{1}\genfrac{}{}{0pt}{}{b}{m}\genfrac{}{}{0pt}{}{c}{n}\bigr)$.

	% page 142

	$$a_0+\frac1{\displaystyle a_1 +
		{\strut \frac1{\displaystyle a_2 +
				{\strut \frac1{\displaystyle a_3 +
						{\strut \frac1{\displaystyle a_4}}}}}}}$$

	% page 143

	$$\binom{p}{2}x^2y^{p-2} - \frac1{1 - x}\frac{1}{1 - x^2}
	=
	\frac{a+1}{b}\bigg/\frac{c+1}{d}.$$

	%% page 145

	$$\sqrt{1+\sqrt{1+\sqrt{1+\sqrt{1+\sqrt{1+x}}}}}$$

	$$\sqrt[n]{1+\sqrt[k]{1+\sqrt[5]{1+\sqrt[4]{1+\sqrt[3]{1+x}}}}}$$

	%% page 147

	$$\left(\frac{\partial^2}{\partial x^2} + \frac{\partial^2}{\partial y^2}\right)
	\bigl|\varphi(x+\mathup{i}y)\bigr|^2=0$$

	%% page 149

	% $$\pi(n)=\sum_{m=2}^n\left\lfloor\biggl(\sum_{k=1}^{m-1}\bigl
	% \lfloor(m/k)\big/\lceil m/k\rceil\bigr\rfloor\biggr)^{-1}\right\rfloor.$$

	$$\pi(n)=\sum_{m=2}^n\left\lfloor\Biggl(\sum_{k=1}^{m-1}\bigl
	\lfloor(m \divslash k) \big/ \lceil m \divslash k\rceil\bigr\rfloor\Biggr)^{-1}\right\rfloor.$$

	% page 168

	$$\int_0^\infty \frac{t - \mathup{i} b}{t^2 + b^2}e^{\mathup{i}at}\,\mathup{d}t=e^{ab}E_1(ab), \quad
	a,b > 0.$$

	% page 176

	$$\mathbf{A} \coloneq \begin{pmatrix}x-\lambda&1&0\\
	0&x-\lambda&1\\
	0&0&x-\lambda\end{pmatrix}.$$

	$$\left\lgroup\begin{matrix}a&b&c\\ d&e&f\\\end{matrix}\right\rgroup
	\left\lgroup\begin{matrix}u&x\cr v&y\cr w&z\end{matrix}\right\rgroup$$

	% page 177

	$$\mathbf{A} = \begin{pmatrix}a_{11}&a_{12}&\ldots&a_{1n}\\
	a_{21}&a_{22}&\ldots&a_{2n}\\
	\vdots&\vdots&\ddots&\vdots\\
	a_{m1}&a_{m2}&\ldots&a_{mn}\end{pmatrix}$$

	$$\mathbf{M}=\bordermatrix{&C&I&C'\cr
		C&1&0&0\cr I&b&1-b&0\cr C'&0&a&1-a}$$

	%% page 186

	$$\sum_{n=0}^\infty a_nz^n\quad\hbox{converges if}\quad
	|z|<\Bigl(\limsup_{n\to\infty}\root n\of{|a_n|}\,\Bigr)^{-1}.$$

	$$\frac{f(x+\mathup{\Delta} x)-f(x)}{\mathup{\Delta} x}\to f'(x)
	\qquad \hbox{as $\mathup{\Delta} x\to0$.}$$

	$$\|u_i\|=1,\qquad u_i\cdot u_j=0\quad\hbox{if $i\ne j$.}$$

	%% page 191

	$$\hbox{The confluent image of}\quad
	\begin{Bmatrix}\hbox{an arc}\hfill\\%
	\hbox{a circle}\hfill\\%
	\hbox{a fan}\hfill\end{Bmatrix}
	\quad\hbox{is}\quad
	\begin{Bmatrix}\hbox{an arc}\hfill\\%
	\hbox{an arc or a circle}\hfill\\%
	\hbox{a fan or an arc}\hfill\end{Bmatrix}.$$

	%% page 191

	\begin{align*}
	T(n)\le T(2^{\lceil\lg n\rceil})
	&\le c(3^{\lceil\lg n\rceil}-2^{\lceil\lg n\rceil})\\
	&<3c\cdot3^{\lg n}\\
	&=3c\,n^{\lg3}.
	\end{align*}

	%\begin{align*}
	%\left\{%
	%\begin{gathered}\alpha&=f(z)\\ \beta&=f(z^2)\\ \gamma&=f(z^3)
	%\end{gathered}
	%\right\}
	%\qquad
	%\left\{%
	%\begin{gathered}
	%x&=\alpha^2-\beta\\ y&=2\gamma
	%\end{gathered}
	%\right\}%
	%\end{align*}

	%$$\left\{
	%\begin{align}
	%\alpha&=f(z)\cr \beta&=f(z^2)\cr \gamma&=f(z^3)\\
	%%\end{align}
	%\right\}
	%\qquad
	%\left\{
	%%\begin{align}
	%x&=\alpha^2-\beta\cr y&=2\gamma\\
	%\end{align}
	%\right\}.$$
	%%% page 192

	\begin{align*}
	\begin{aligned}
	(x+y)(x-y)&=x^2-xy+yx-y^2\\
	&=x^2-y^2\\
	(x+y)^2&=x^2+2xy+y^2.
	\end{aligned}
	\end{align*}

	%% page 192

	\begin{align*}
	\begin{aligned}
	\left( \int\limits_{-\infty}^\infty \mathup{e}^{-x^2}\,\mathup{d}x \right)^2
	&=\int_{-\infty}^\infty\int_{-\infty}^\infty \mathup{e}^{-(x^2+y^2)}\,\mathup{d}x\,\mathup{d}y\\
	&=\int_0^{2\piup}\int_0^\infty \mathup{e}^{-r^2}\,\mathup{d}r\,\mathup{d}\theta\\
	&=\int_0^{2\piup}\biggl(\mathup{e}^{-\frac{r^2}{2}}\biggl|_{r=0}^{r=\infty}\,\biggr)\,\mathup{d}\theta\\
	&=\piup.
	\end{aligned}
	\end{align*}


	%% page 197

	$$\prod_{k\ge0}\frac{1}{(1-q^kz)}=
	\sum_{n\ge0}z^n \bigg/ \!\!\prod_{1\le k\le n}(1-q^k).$$

	$$\sum_{\substack{\scriptstyle 0< i\le m\\\scriptstyle0<j\le n}}p(i,j) \,\ne
	%
	% $$\sum_{i=1}^p \sum_{j=1}^q \sum_{k=1}^r a_{ij} b_{jk} c_{ki}$$
	%
	\sum_{i=1}^p \sum_{j=1}^q \sum_{k=1}^r a_{ij} b_{jk} c_{ki} \,\ne
	%
	\sum_{\substack{\scriptstyle 1\le i\le p \\ \scriptstyle 1\le j\le q\\
			\scriptstyle 1\le k\le r}} a_{ij} b_{jk} c_{ki}$$

	$$\max_{1\le n\le m}\log_2P_n \quad \hbox{and} \quad
	\lim_{x\to0}\frac{\sin x}{x}=1$$
	Inline math:
	$\max_{1\le n\le m}\log_2P_n \quad \hbox{and} \quad
	\lim_{x\to0}\frac{\sin x}{x}=1$
	$$p_1(n)=\lim_{m\to\infty}\sum_{\nu=0}^\infty\bigl(1-\cos^{2m}(\nu!^n\piup \divslash n)\bigr)$$
	Inline math:
	$p_1(n)=\lim_{m\to\infty}\sum_{\nu=0}^\infty\bigl(1-\cos^{2m}(\nu!^n\piup \divslash n)\bigr)$

	\clearpage

	\begingroup
	\bfseries\boldmath

	\section{Math \texorpdfstring{``}{“}Torture Test\texorpdfstring{''}{”} \texttt{\textbackslash boldmath}}

	Most of the following examples are taken from \textit{The \TeX book} \citep[][see \url{https://ctan.org/pkg/texbook}]{knuth:ct:a} and were adapted for \LaTeX\ from Karl Berry's torture test for plain \TeX\ math fonts.

	\noindent $x + y - z$, \quad $x + y * z$, \quad $z * y \divslash z$, \quad
	$(x+y)(x-y) = x^2 - y^2$,

	\noindent $x \times y \cdot z = [x\, y\, z]$, \quad $x\circ y \bullet z$, \quad
	$x\cup y \cap z$, \quad $x\sqcup y \sqcap z$, \quad

	\noindent $x \vee y \wedge z$, \quad $x\pm y\mp z$, \quad
	$x=y \divslash z$, \;\; $x:=y$, \;\; $x\le y \ne z$, \;\; $x \sim y \simeq z$
	$x \equiv y \nequiv z$, \;\; $x\subset y \subseteq z$

	\noindent $\sin2\theta=2\sin\theta\cos\theta$, \quad
	$\hbox{O}(n\log n\log n)$, \quad
	$\Pr(X>x)=\exp(-x \divslash \mu)$,

	\noindent $\bigl(x\in A(n)\bigm|x\in B(n)\bigr)$, \quad
	$\bigcup_n X_n\bigm\|\bigcap_n Y_n$

	% page 178

	\noindent In-text matrices $\binom{1\,1}{0\,1}$ and $\bigl(\genfrac{}{}{0pt}{}{a}{1}\genfrac{}{}{0pt}{}{b}{m}\genfrac{}{}{0pt}{}{c}{n}\bigr)$.

	% page 142

	$$a_0+\frac1{\displaystyle a_1 +
		{\strut \frac1{\displaystyle a_2 +
				{\strut \frac1{\displaystyle a_3 +
						{\strut \frac1{\displaystyle a_4}}}}}}}$$

	% page 143

	$$\binom{p}{2}x^2y^{p-2} - \frac1{1 - x}\frac{1}{1 - x^2}
	=
	\frac{a+1}{b}\bigg/\frac{c+1}{d}.$$

	%% page 145

	$$\sqrt{1+\sqrt{1+\sqrt{1+\sqrt{1+\sqrt{1+x}}}}}$$

	$$\sqrt[n]{1+\sqrt[k]{1+\sqrt[5]{1+\sqrt[4]{1+\sqrt[3]{1+x}}}}}$$

	%% page 147

	$$\left(\frac{\partial^2}{\partial x^2} + \frac{\partial^2}{\partial y^2}\right)
	\bigl|\varphi(x+\mathup{i}y)\bigr|^2=0$$

	%% page 149

	% $$\pi(n)=\sum_{m=2}^n\left\lfloor\biggl(\sum_{k=1}^{m-1}\bigl
	% \lfloor(m/k)\big/\lceil m/k\rceil\bigr\rfloor\biggr)^{-1}\right\rfloor.$$

	$$\pi(n)=\sum_{m=2}^n\left\lfloor\Biggl(\sum_{k=1}^{m-1}\bigl
	\lfloor(m \divslash k) \big/ \lceil m \divslash k\rceil\bigr\rfloor\Biggr)^{-1}\right\rfloor.$$

	% page 168

	$$\int_0^\infty \frac{t - \mathup{i} b}{t^2 + b^2}e^{\mathup{i}at}\,\mathup{d}t=e^{ab}E_1(ab), \quad
	a,b > 0.$$

	% page 176

	$$\mathbf{A} \coloneq \begin{pmatrix}x-\lambda&1&0\\
	0&x-\lambda&1\\
	0&0&x-\lambda\end{pmatrix}.$$

	$$\left\lgroup\begin{matrix}a&b&c\\ d&e&f\\\end{matrix}\right\rgroup
	\left\lgroup\begin{matrix}u&x\cr v&y\cr w&z\end{matrix}\right\rgroup$$

	% page 177

	$$\mathbf{A} = \begin{pmatrix}a_{11}&a_{12}&\ldots&a_{1n}\\
	a_{21}&a_{22}&\ldots&a_{2n}\\
	\vdots&\vdots&\ddots&\vdots\\
	a_{m1}&a_{m2}&\ldots&a_{mn}\end{pmatrix}$$

	$$\mathbf{M}=\bordermatrix{&C&I&C'\cr
		C&1&0&0\cr I&b&1-b&0\cr C'&0&a&1-a}$$

	%% page 186

	$$\sum_{n=0}^\infty a_nz^n\quad\hbox{converges if}\quad
	|z|<\Bigl(\limsup_{n\to\infty}\root n\of{|a_n|}\,\Bigr)^{-1}.$$

	$$\frac{f(x+\mathup{\Delta} x)-f(x)}{\mathup{\Delta} x}\to f'(x)
	\qquad \hbox{as $\mathup{\Delta} x\to0$.}$$

	$$\|u_i\|=1,\qquad u_i\cdot u_j=0\quad\hbox{if $i\ne j$.}$$

	%% page 191

	$$\hbox{The confluent image of}\quad
	\begin{Bmatrix}\hbox{an arc}\hfill\\\hbox{a circle}\hfill\\
	\hbox{a fan}\hfill\\\end{Bmatrix}
	\quad\hbox{is}\quad
	\begin{Bmatrix}\hbox{an arc}\hfill\\
	\hbox{an arc or a circle}\hfill\\
	\hbox{a fan or an arc}\hfill\end{Bmatrix}.$$

	%% page 191

	\begin{align*}
	T(n)\le T(2^{\lceil\lg n\rceil})
	&\le c(3^{\lceil\lg n\rceil}-2^{\lceil\lg n\rceil})\\
	&<3c\cdot3^{\lg n}\\
	&=3c\,n^{\lg3}.
	\end{align*}

	%\begin{align*}
	%\left\{%
	%\begin{gathered}\alpha&=f(z)\\ \beta&=f(z^2)\\ \gamma&=f(z^3)
	%\end{gathered}
	%\right\}
	%\qquad
	%\left\{%
	%\begin{gathered}
	%x&=\alpha^2-\beta\\ y&=2\gamma
	%\end{gathered}
	%\right\}%
	%\end{align*}

	%$$\left\{
	%\begin{align}
	%\alpha&=f(z)\cr \beta&=f(z^2)\cr \gamma&=f(z^3)\\
	%%\end{align}
	%\right\}
	%\qquad
	%\left\{
	%%\begin{align}
	%x&=\alpha^2-\beta\cr y&=2\gamma\\
	%\end{align}
	%\right\}.$$
	%%% page 192

	\begin{align*}
	\begin{aligned}
	(x+y)(x-y)&=x^2-xy+yx-y^2\\
	&=x^2-y^2\\
	(x+y)^2&=x^2+2xy+y^2.
	\end{aligned}
	\end{align*}

	%% page 192

	\begin{align*}
	\begin{aligned}
	\left( \int\limits_{-\infty}^\infty \mathup{e}^{-x^2}\,\mathup{d}x \right)^2
	&=\int_{-\infty}^\infty\int_{-\infty}^\infty \mathup{e}^{-(x^2+y^2)}\,\mathup{d}x\,\mathup{d}y\\
	&=\int_0^{2\piup}\int_0^\infty \mathup{e}^{-r^2}\,\mathup{d}r\,\mathup{d}\theta\\
	&=\int_0^{2\piup}\biggl(\mathup{e}^{-\frac{r^2}{2}}\biggl|_{r=0}^{r=\infty}\,\biggr)\,\mathup{d}\theta\\
	&=\piup.
	\end{aligned}
	\end{align*}

	%% page 197

	$$\prod_{k\ge0}\frac{1}{(1-q^kz)}=
	\sum_{n\ge0}z^n \bigg/ \!\!\prod_{1\le k\le n}(1-q^k).$$

	$$\sum_{\substack{\scriptstyle 0< i\le m\\\scriptstyle0<j\le n}}p(i,j) \,\ne
	%
	% $$\sum_{i=1}^p \sum_{j=1}^q \sum_{k=1}^r a_{ij} b_{jk} c_{ki}$$
	%
	\sum_{i=1}^p \sum_{j=1}^q \sum_{k=1}^r a_{ij} b_{jk} c_{ki} \,\ne
	%
	\sum_{\substack{\scriptstyle 1\le i\le p \\ \scriptstyle 1\le j\le q\\
			\scriptstyle 1\le k\le r}} a_{ij} b_{jk} c_{ki}$$

	\medskip
	$$\max_{1\le n\le m}\log_2P_n \quad \hbox{and} \quad
	\lim_{x\to0}\frac{\sin x}{x}=1$$

	\medskip
	Inline math:
	$\max_{1\le n\le m}\log_2P_n \quad \hbox{and} \quad
	\lim_{x\to0}\frac{\sin x}{x}=1$
	$$p_1(n)=\lim_{m\to\infty}\sum_{\nu=0}^\infty\bigl(1-\cos^{2m}(\nu!^n\piup \divslash n)\bigr)$$

	\medskip
	Inline math:
	$p_1(n)=\lim_{m\to\infty}\sum_{\nu=0}^\infty\bigl(1-\cos^{2m}(\nu!^n\piup \divslash n)\bigr)$

	\endgroup

	\clearpage

	\section{Extensive Math Test \showfamily}

	\subsection{Spacing}

	$$\frac{a \divslash b + \frac{a \divslash b + c}{x}}{x} \qquad \sin x \divslash \cos x \qquad n \divslash \log n$$

	{\boldmath $$\frac{a \divslash b + \frac{a \divslash b + c}{x}}{x} \qquad \sin x \divslash \cos x \qquad n \divslash \log n$$}

	\begin{theorem}[simplest form of the Central Limit Theorem]
		\label{theorem:SFCLT}
		Let $X_1, X_2, \ldots, X_n$ be a~sequence of i.i.d. random variables with mean $0$
		and variance $1$ on a~probability space $(\Omega, \mathcal{F}, \mathbb{P})$. Then
		\[
			\mathbb{P}{\left(\frac{X_1+\cdots+X_n}{\sqrt{n}}\le y\right)} \to
			\mathfrak{N}(y) \coloneq
			\int_{-\infty}^y \frac{\mathup{e}^{-v^2 \divslash 2}}{\sqrt{2\mathup{\pi}}}\,\mathup{d}v
			\quad\text{as} \quad n\to\infty,
		\]
		or, equivalently, letting $S_n \coloneq \sum_1^n X_k$,
		\[
			\mathbb{E} f \big( S_n \divslash \sqrt{n} \big) \to
			\int_{-\infty}^\infty f(v) \frac{\mathup{e}^{-v^2 \divslash 2}}{\sqrt{2\piup}}\,\mathup{d}v
			\quad \mbox{as $n\to\infty$, for every $f \in \mathup{b}\mathcal{C}(\mathbb{R})$.}
		\]
	\end{theorem}

	\subsection{Formulas \showfamily}

	\setlength{\parindent}{0pt}

	\noindent%
	\checkgreekletters

	\noindent%
	{\boldmath\checkgreekletters}

	\noindent%
	{\checkgreeklettersup}

	\noindent%
	{\boldmath \checkgreeklettersup}

	\noindent%
	$\alpha a > 0, \beta b + (3 \times 27), \Gamma G = 7 < 8, \lambda$

	$\lim_{\nu \to \infty} v(\nu) = \max_{s \in S} \{s \pm 3 \gamma + y - 1\} = 4 \times 7$

	$\hat{\bm{\beta}} = (\bm{X}'\bm{X})^{-1}\bm{X}'\bm{y}$

	$$\lim_{N \to \infty} \sum_{i=0}^{N} x^i = \min_{x \in \mathbb{R}} S(x)$$

	$$\int_{-\infty}^{\infty} x\,f(x)\,\mathup{d}x = \left( \frac{27}{2} \right)$$

	$$t[u_1, \dots, u_n] = \sum_{k=1}^n \binom{n - 1}{k - 1} \, (1 - t)^{n-k} \, t^{k-1} \, u_k$$

	$$\Phi(u) = \frac{1}{\sqrt{2\uppi}} \int_{-\infty}^u \mathrm{e}^{-t^2 \divslash 2} \, \dd t$$

	Disambiguation: $0$~O~$O$, $1$~l~I~$|$~$l$~$I$~$/$, $i$~$j$, $rn$~$m$, $\theta$~$\Theta$, $\phi$~$\psi$, --~$-$

	Latin vs. Greek: $a$~$\alpha$, $d$~$\delta$, $e$~$\epsilon$, $i$~$\iota$, $k$~$\kappa$, $n$~$\eta$, $o$~$\sigma$, $p$~$\rho$, \textit{\ss} $\beta$, $u$~$\upsilon$, $v$~$\nu$, $w$~$\omega$, $x$~$\chi$, $y$~$\gamma$, $A$~$\Delta$~$\Lambda$, $O$~$\Theta$~$\Omega$, $T$~$\Gamma$, $Y$~$\Upsilon$.

	\noindent%
	{\bfseries\boldmath%
		$\alpha a > 0, \beta b + (3 \times 27), \Gamma G = 7 < 8, \lambda$

		$\lim_{\nu \to \infty} v(\nu) = \max_{s \in S} \{s \pm 3 \gamma + y - 1\} = 4 \times 7$

		$\hat{\bm{\beta}} = (\bm{X}'\bm{X})^{-1}\bm{X}'\bm{y}$

		$$\lim_{N \to \infty} \sum_{i=0}^{N} x^i = \min_{x \in \mathbb{R}} S(x)$$

		$$\int_{-\infty}^{\infty} x\,f(x)\,\mathup{d}x = \left( \frac{27}{2} \right)$$

		$$t[u_1, \dots, u_n] = \sum_{k=1}^n \binom{n - 1}{k - 1} \, (1 - t)^{n-k} \, t^{k-1} \, u_k$$

		$$\Phi(u) = \frac{1}{\sqrt{2\uppi}} \int_{-\infty}^u \mathrm{e}^{-t^2 \divslash 2} \, \dd t$$

		Disambiguation: $0$~O~$O$, $1$~l~I~$|$~$l$~$I$~$/$, $i$~$j$, $rn$~$m$, $\theta$~$\Theta$, $\phi$~$\psi$, --~$-$

		Latin vs. Greek: $a$~$\alpha$, $d$~$\delta$, $e$~$\epsilon$, $i$~$\iota$, $k$~$\kappa$, $n$~$\eta$, $o$~$\sigma$, $p$~$\rho$, \textit{\ss} $\beta$, $u$~$\upsilon$, $v$~$\nu$, $w$~$\omega$, $x$~$\chi$, $y$~$\gamma$, $A$~$\Delta$~$\Lambda$, $O$~$\Theta$~$\Omega$, $T$~$\Gamma$, $Y$~$\Upsilon$.
	}

	\subsection{Math Alphabets \showfamily}
	\label{app:math-test:alphabets}

	%\sffamily\selectfont

	Default
	\def\test#1{#1,}
	\begin{eqnarray*}
	  && {\testnums}\\
	  && {\testupper}\\
	  && {\testlower}\\
	  && {\testupgreek}\\
	  && {\testlowgreek}
	\end{eqnarray*}%

	Math Normal (\texttt{\string\mathnormal})
	\def\test#1{\mathnormal{#1},}
	\begin{eqnarray*}
	  && {\testnums}\\
	  && {\testupper}\\
	  && {\testlower}\\
	  && {\testupgreek}\\
	  && {\testlowgreek}
	\end{eqnarray*}%

	Math Italic (\texttt{\string\mathit})
	\def\test#1{\mathit{#1},}
	\begin{eqnarray*}
	  && {\testnums}\\
	  && {\testupper}\\
	  && {\testlower}\\
	  && {\testupgreek}\\
	  && {\testlowgreek}
	\end{eqnarray*}%

	Math Roman (\texttt{\string\mathrm})
	\def\test#1{\mathrm{#1},}
	\begin{eqnarray*}
	  && {\testnums}\\
	  && {\testupper}\\
	  && {\testlower}\\
	  && {\testupgreek}\\
	  && {\testlowgreekup}
	\end{eqnarray*}%

	Math Italic Bold (\texttt{\string\bm})
	\def\test#1{\bm{#1},}
	\begin{eqnarray*}
	  && {\testnums}\\
	  && {\testupper}\\
	  && {\testlower}\\
	  && {\testupgreek}\\
	  && {\testlowgreek}
	\end{eqnarray*}%

	Math Bold (\texttt{\string\mathbf})
	\def\test#1{\mathbf{#1},}
	\begin{eqnarray*}
	  && {\testnums}\\
	  && {\testupper}\\
	  && {\testlower}\\
	  && {\testupgreek}\\
	  && {\def\test#1{\bm{\mathbf{#1}},}\testlowgreekup}
	\end{eqnarray*}%

	Caligraphic (\texttt{\string\mathcal})
	\def\test#1{\mathcal{#1},}
	\begin{eqnarray*}
	  && {\testupper}
	\end{eqnarray*}%

	Script (\texttt{\string\mathscr})
	\def\test#1{\mathscr{#1},}
	\begin{eqnarray*}
	  && {\testupper}
	\end{eqnarray*}%

	Fraktur (\texttt{\string\mathfrak})
	\def\test#1{\mathfrak{#1},}
	\begin{eqnarray*}
	  && {\testupper}\\
	  && {\testlower}
	\end{eqnarray*}%

	Blackboard Bold (\texttt{\string\mathbb})
	\def\test#1{\mathbb{#1},}
	\begin{eqnarray*}
	  && {\testupper}
	\end{eqnarray*}%

	\subsection{Character Sidebearings \showfamily}

	Default
	\def\test#1{|#1|+{}}
	\begin{eqnarray*}
	  && {\testupperi}\\
	  && {\testupperii}\\
	  && {\testloweri}\\
	  && {\testlowerii}\\
	  && {\testupgreeki}\\
	  && {\testupgreekii}\\
	  && {\testlowgreeki}\\
	  && {\testlowgreekii}\\
	  && {\testlowgreekiii}
	\end{eqnarray*}%

	Math Roman (\texttt{\string\mathrm})
	\def\test#1{|\mathrm{#1}|+{}}%
	\begin{eqnarray*}
	  && {\testupperi}\\
	  && {\testupperii}\\
	  && {\testloweri}\\
	  && {\testlowerii}\\
	  && {\testupgreeki}\\
	  && {\testupgreekii}\\
	  && {\testlowgreekupi}\\
	  && {\testlowgreekupii}\\
	  && {\testlowgreekupiii}
	\end{eqnarray*}%

	Math Italic Bold (\texttt{\string\bm})
	\def\test#1{|\bm{#1}|+{}}%
	\begin{eqnarray*}
	  && {\testupperi}\\
	  && {\testupperii}\\
	  && {\testloweri}\\
	  && {\testlowerii}\\
	  && {\testupgreeki}\\
	  && {\testupgreekii}\\
	  && {\testlowgreeki}\\
	  && {\testlowgreekii}\\
	  && {\testlowgreekiii}
	\end{eqnarray*}%

	Math Bold (\texttt{\string\mathbf})
	\def\test#1{|\mathbf{#1}|+{}}%
	\begin{eqnarray*}
	  && {\testupperi}\\
	  && {\testupperii}\\
	  && {\testloweri}\\
	  && {\testlowerii}\\
	  && {\testupgreeki}\\
	  && {\testupgreekii}\\
	  && {\def\test#1{|\bm{\mathbf{#1}}|+{}}\testlowgreekupi}\\
	  && {\def\test#1{|\bm{\mathbf{#1}}|+{}}\testlowgreekupii}\\
	  && {\def\test#1{|\bm{\mathbf{#1}}|+{}}\testlowgreekupiii}
	\end{eqnarray*}%

	Math Calligraphic (\texttt{\string\mathcal})
	\def\test#1{|\mathcal{#1}|+{}}%
	\begin{eqnarray*}
	  && {\testupperi}\\
	  && {\testupperii}
	\end{eqnarray*}%

	\subsection{Superscript Positioning \showfamily}

	Default
	\def\test#1{#1^{2}+{}}%
	\begin{eqnarray*}
	  && {\testupperi}\\
	  && {\testupperii}\\
	  && {\testloweri}\\
	  && {\testlowerii}\\
	  && {\testupgreeki}\\
	  && {\testupgreekii}\\
	  && {\testlowgreeki}\\
	  && {\testlowgreekii}\\
	  && {\testlowgreekiii}
	\end{eqnarray*}%

	Math Roman (\texttt{\string\mathrm})
	\def\test#1{\mathrm{#1}^{2}+{}}%
	\begin{eqnarray*}
	  && {\testupperi}\\
	  && {\testupperii}\\
	  && {\testloweri}\\
	  && {\testlowerii}\\
	  && {\testupgreeki}\\
	  && {\testupgreekii}\\
	  && {\testlowgreekupi}\\
	  && {\testlowgreekupii}\\
	  && {\testlowgreekupiii}
	\end{eqnarray*}%

	Math Italic Bold (\texttt{\string\bm})
	\def\test#1{\bm{#1}^{2}+{}}%
	\begin{eqnarray*}
	  && {\testupperi}\\
	  && {\testupperii}\\
	  && {\testloweri}\\
	  && {\testlowerii}\\
	  && {\testupgreeki}\\
	  && {\testupgreekii}\\
	  && {\testlowgreeki}\\
	  && {\testlowgreekii}\\
	  && {\testlowgreekiii}
	\end{eqnarray*}%

	Math Bold (\texttt{\string\mathbf})
	\def\test#1{\mathbf{#1}^{2}+{}}%
	\begin{eqnarray*}
	  && {\testupperi}\\
	  && {\testupperii}\\
	  && {\testloweri}\\
	  && {\testlowerii}\\
	  && {\testupgreeki}\\
	  && {\testupgreekii}\\
	  && {\def\test#1{\bm{\mathbf{#1}}^{2}+{}}\testlowgreekupi}\\
	  && {\def\test#1{\bm{\mathbf{#1}}^{2}+{}}\testlowgreekupii}\\
	  && {\def\test#1{\bm{\mathbf{#1}}^{2}+{}}\testlowgreekupiii}
	\end{eqnarray*}

	Math Calligraphic (\texttt{\string\mathcal})
	\def\test#1{\mathcal{#1}^{2}+{}}%
	\begin{eqnarray*}
	  && {\testupperi}\\
	  && {\testupperii}
	\end{eqnarray*}%

	\subsection{Subscript Positioning \showfamily}

	Default
	\def\test#1{\mathnormal{#1}_{i}+{}}%
	\begin{eqnarray*}
	  && {\testupperi}\\
	  && {\testupperii}\\
	  && {\testloweri}\\
	  && {\testlowerii}\\
	  && {\testupgreeki}\\
	  && {\testupgreekii}\\
	  && {\testlowgreeki}\\
	  && {\testlowgreekii}\\
	  && {\testlowgreekiii}
	\end{eqnarray*}%

	Math Roman (\texttt{\string\mathrm})
	\def\test#1{\mathrm{#1}_{i}+{}}%
	\begin{eqnarray*}
	  && {\testupperi}\\
	  && {\testupperii}\\
	  && {\testloweri}\\
	  && {\testlowerii}\\
	  && {\testupgreeki}\\
	  && {\testupgreekii}\\
	  && {\testlowgreekupi}\\
	  && {\testlowgreekupii}\\
	  && {\testlowgreekupiii}
	\end{eqnarray*}%

	Math Bold Italic (\texttt{\string\bm})
	\def\test#1{\bm{#1}_{i}+{}}%
	\begin{eqnarray*}
	  && {\testupperi}\\
	  && {\testupperii}\\
	  && {\testloweri}\\
	  && {\testlowerii}\\
	  && {\testupgreeki}\\
	  && {\testupgreekii}\\
	  && {\testlowgreeki}\\
	  && {\testlowgreekii}\\
	  && {\testlowgreekiii}
	\end{eqnarray*}

	Math Bold (\texttt{\string\mathbf})
	\def\test#1{\mathbf{#1}_{i}+{}}%
	\begin{eqnarray*}
	  && {\testupperi}\\
	  && {\testupperii}\\
	  && {\testloweri}\\
	  && {\testlowerii}\\
	  && {\testupgreeki}\\
	  && {\testupgreekii}\\
	  && {\def\test#1{\bm{\mathbf{#1}}_{i}+{}}\testlowgreekupi}\\
	  && {\def\test#1{\bm{\mathbf{#1}}_{i}+{}}\testlowgreekupii}\\
	  && {\def\test#1{\bm{\mathbf{#1}}_{i}+{}}\testlowgreekupiii}
	\end{eqnarray*}%

	Math Calligraphic (\texttt{\string\mathcal})
	\def\test#1{\mathcal{#1}_{i}+{}}%
	\begin{eqnarray*}
	  && {\testupperi}\\
	  && {\testupperii}
	\end{eqnarray*}%

	\subsection{Accent Positioning \showfamily}

	Default
	\def\test#1{\hat{#1}+{}}%
	\begin{eqnarray*}
	  && {\testnums}\\
	  && {\testupperi}\\
	  && {\testupperii}\\
	  && {\testloweri}\\
	  && {\testlowerii}\\
	  && {\testupgreeki}\\
	  && {\testupgreekii}\\
	  && {\testlowgreeki}\\
	  && {\testlowgreekii}\\
	  && {\testlowgreekiii}
	\end{eqnarray*}%

	Math Italic (\texttt{\string\mathit})
	\def\test#1{\hat{\mathit{#1}}+{}}%
	\begin{eqnarray*}
	  && {\testnums}\\
	  && {\testupperi}\\
	  && {\testupperii}\\
	  && {\testloweri} \test\ell \test\wp \test\imath \test\jmath \tilde{i} \\
	  && {\testlowerii}\\
	  && {\testupgreeki}\\
	  && {\testupgreekii}\\
	  && {\testlowgreeki}\\
	  && {\testlowgreekii}\\
	  && {\testlowgreekiii}
	\end{eqnarray*}%

	Math Roman (\texttt{\string\mathrm})
	\def\test#1{\hat{\mathrm{#1}}+{}}%
	\begin{eqnarray*}
	  && {\testnums}\\
	  && {\testupperi}\\
	  && {\testupperii}\\
	  && {\testloweri}\\
	  && {\testlowerii}\\
	  && {\testupgreeki}\\
	  && {\testupgreekii}\\
	  && {\testlowgreekupi}\\
	  && {\testlowgreekupii}\\
	  && {\testlowgreekupiii}
	\end{eqnarray*}%

	Math Italic Bold (\texttt{\string\bm})
	\def\test#1{\hat{\bm{#1}}+{}}%
	\begin{eqnarray*}
	  && {\testnums}\\
	  && {\testupperi}\\
	  && {\testupperii}\\
	  && {\testloweri}\\
	  && {\testlowerii}\\
	  && {\testupgreeki}\\
	  && {\testupgreekii}\\
	  && {\testlowgreeki}\\
	  && {\testlowgreekii}\\
	  && {\testlowgreekiii}
	\end{eqnarray*}%

	Math Bold (\texttt{\string\mathbf})
	\def\test#1{\hat{\mathbf{#1}}+{}}%
	\begin{eqnarray*}
	  && {\testnums}\\
	  && {\testupperi}\\
	  && {\testupperii}\\
	  && {\testloweri}\\
	  && {\testlowerii}\\
	  && {\testupgreeki}\\
	  && {\testupgreekii}\\
	  && {\def\test#1{\hat{\bm{\mathbf{#1}}}+{}}\testlowgreekupi}\\
	  && {\def\test#1{\hat{\bm{\mathbf{#1}}}+{}}\testlowgreekupii}\\
	  && {\def\test#1{\hat{\bm{\mathbf{#1}}}+{}}\testlowgreekupiii}
	\end{eqnarray*}

	Math Calligraphic (\texttt{\string\mathcal})
	\def\test#1{\hat{\mathcal{#1}}+{}}%
	\begin{eqnarray*}
	  && {\testupperi}\\
	  && {\testupperii}
	\end{eqnarray*}%

	\subsection{Differentials \showfamily}

	\begin{eqnarray*}
	\gdef\test#1{\dd #1+{}}%
	  && {\testupperi}\\
	  && {\testupperii}\\
	  && {\testloweri}\\
	  && {\testlowerii}\\
	  && {\testupgreeki}\\
	  && {\testupgreekii}\\
	  && {\testlowgreeki}\\
	  && {\testlowgreekii}\\
	  && {\testlowgreekiii}\\
	\gdef\test#1{\dd \mathrm{#1}+{}}%
	  && {\testupgreeki}\\
	  && {\testupgreekii}\\
	  && {\testlowgreekupi}\\
	  && {\testlowgreekupii}\\
	  && {\testlowgreekupiii}\\
	\end{eqnarray*}%

	\begin{eqnarray*}
	\gdef\test#1{\partial #1+{}}%
	  && {\testupperi}\\
	  && {\testupperii}\\
	  && {\testloweri}\\
	  && {\testlowerii}\\
	  && {\testupgreeki}\\
	  && {\testupgreekii}\\
	  && {\testlowgreeki}\\
	  && {\testlowgreekii}\\
	  && {\testlowgreekiii}\\
	\gdef\test#1{\partial \mathrm{#1}+{}}%
	  && {\testupgreeki}\\
	  && {\testupgreekii}\\
	  && {\testlowgreekupi}\\
	  && {\testlowgreekupii}\\
	  && {\testlowgreekupiii}\\
	\end{eqnarray*}%

	\subsection{Slash Kerning \showfamily}

	\def\test#1{1 \divslash #1+{}}
	\begin{eqnarray*}
	  && {\testupperi}\\
	  && {\testupperii}\\
	  && {\testloweri}\\
	  && {\testlowerii}\\
	  && {\testupgreeki}\\
	  && {\testupgreekii}\\
	  && {\testlowgreeki}\\
	  && {\testlowgreekii}\\
	  && {\testlowgreekiii}
	\end{eqnarray*}

	\def\test#1{#1 \divslash 2+{}}
	\begin{eqnarray*}
	  && {\testupperi}\\
	  && {\testupperii}\\
	  && {\testloweri}\\
	  && {\testlowerii}\\
	  && {\testupgreeki}\\
	  && {\testupgreekii}\\
	  && {\testlowgreeki}\\
	  && {\testlowgreekii}\\
	  && {\testlowgreekiii}
	\end{eqnarray*}

	\subsection{(Big) Operators \showfamily}

	\def\testop#1{#1_{i=1}^{n} x^{n} \quad}
	$
		\testop\sum
		\testop\prod
		\testop\coprod
		\testop\int
		\testop\oint
	$

	\noindent%
	$
		\testop\bigotimes
		\testop\bigoplus
		\testop\bigodot
		\testop\bigwedge
		\testop\bigvee
		\testop\biguplus
		\testop\bigcup
		\testop\bigcap
		\testop\bigsqcup
		% \testop\bigsqcap
	$

	\begin{displaymath}
	  \testop\sum
	  \testop\prod
	  \testop\coprod
	  \testop\int
	  \testop\oint
	\end{displaymath}
	\begin{displaymath}
	  \testop\bigotimes
	  \testop\bigoplus
	  \testop\bigodot
	  \testop\bigwedge
	  \testop\bigvee
	  \testop\biguplus
	  \testop\bigcup
	  \testop\bigcap
	  \testop\bigsqcup
	% \testop\bigsqcap
	\end{displaymath}

	\subsection{Radicals \showfamily}

	\begin{displaymath}
	  \sqrt{x+y} \qquad \sqrt{x^{2}+y^{2}} \qquad
	  \sqrt{x_{i}^{2}+y_{j}^{2}} \qquad
	  \sqrt{\left(\frac{\cos x}{2}\right)} \qquad
	  \sqrt{\left(\frac{\sin x}{2}\right)}
	\end{displaymath}

	\begingroup
	\delimitershortfall-1pt
	\begin{displaymath}
	  \sqrt{\sqrt{\sqrt{\sqrt{\sqrt{\sqrt{\sqrt{x+y}}}}}}}
	\end{displaymath}
	\endgroup % \delimitershortfall

	\subsection{Over- and Underbraces \showfamily}

	\begin{displaymath}
	  \overbrace{x} \quad
	  \overbrace{x+y} \quad
	  \overbrace{x^{2}+y^{2}} \quad
	  \overbrace{x_{i}^{2}+y_{j}^{2}} \quad
	  \underbrace{x} \quad
	  \underbrace{x+y} \quad
	  \underbrace{x_{i}+y_{j}} \quad
	  \underbrace{x_{i}^{2}+y_{j}^{2}} \quad
	\end{displaymath}

	\subsection{Normal and Wide Accents \showfamily}

	\begin{displaymath}
	  \dot{x} \quad
	  \ddot{x} \quad
	  \vec{x} \quad
	  \bar{x} \quad
	  \overline{x} \quad
	  \overline{xx} \quad
	  \tilde{x} \quad
	  \widetilde{x} \quad
	  \widetilde{xx} \quad
	  \widetilde{xxx} \quad
	  \hat{x} \quad
	  \widehat{x} \quad
	  \widehat{xx} \quad
	  \widehat{xxx} \quad
	\end{displaymath}

	\begin{displaymath}
	  \hat{x} \quad
	  \check{x} \quad
	  \tilde{x} \quad
	  \acute{x} \quad
	  \grave{x} \quad
	  \dot{x} \quad
	  \ddot{x} \quad
	  \breve{x} \quad
	  \bar{x} \quad
	  \vec{x} \quad
	\end{displaymath}

	\subsection{Long Arrows \showfamily}

	\begin{displaymath}
	  \leftarrow \mathrel{-} \rightarrow \quad
	  \leftrightarrow \quad
	  \longleftarrow  \quad
	  \longrightarrow \quad
	  \longleftrightarrow \quad
	  \Leftarrow = \Rightarrow \quad
	  \Leftrightarrow \quad
	  \Longleftarrow  \quad
	  \Longrightarrow \quad
	  \Longleftrightarrow \quad
	\end{displaymath}

	\subsection{Left and Right Delimiters \showfamily}

	\def\testdelim#1#2{ - #1 f #2 - }
	\begin{displaymath}
	  \testdelim()
	  \testdelim[]
	  \testdelim\lfloor\rfloor
	  \testdelim\lceil\rceil
	  \testdelim\langle\rangle
	  \testdelim\{\}
	\end{displaymath}

	Using {\tt\string\left} and {\tt\string\right}.
	\def\testdelim#1#2{ - \left#1 f \right#2 - }
	\begin{displaymath}
	  \testdelim()
	  \testdelim[]
	  \testdelim\lfloor\rfloor
	  \testdelim\lceil\rceil
	  \testdelim\langle\rangle
	  \testdelim\{\}
	% \testdelim\lgroup\rgroup
	% \testdelim\lmoustache\rmoustache
	\end{displaymath}
	\begin{displaymath}
	  \testdelim)(
	  \testdelim][
	  \testdelim//
	  \testdelim\backslash\backslash
	  \testdelim/\backslash
	  \testdelim\backslash/
	\end{displaymath}

	\subsection{Big-g-g Delimiters \showfamily}

	\def\testdelim#1#2{%
	  - \left#1\left#1\left#1\left#1\left#1\left#1\left#1\left#1 -
	  \right#2\right#2\right#2\right#2\right#2\right#2\right#2\right#2 -}

	\begingroup
	\delimitershortfall-1pt
	\begin{displaymath}
	  \testdelim\lfloor\rfloor
	  \qquad
	  \testdelim()
	\end{displaymath}
	\begin{displaymath}
	  \testdelim\lceil\rceil
	  \qquad
	  \testdelim\{\}
	\end{displaymath}
	\begin{displaymath}
	  \testdelim[]
	  \qquad
	  \testdelim\lgroup\rgroup
	\end{displaymath}
	\begin{displaymath}
	  \testdelim\langle\rangle
	  \qquad
	  \testdelim\lmoustache\rmoustache
	\end{displaymath}
	\begin{displaymath}
	  \testdelim\uparrow\downarrow \quad
	  \testdelim\Uparrow\Downarrow \quad
	\end{displaymath}
	\endgroup % \delimitershortfall

	\def\X#1{$x #1 y$ &\tt\string#1}
	\def\Y#1{$\big#1$ &\tt\string#1}
	\def\Z#1{$x #1 y$}
	\def\W#1#2{$#1{#2}$ &\tt\string#1\string{#2\string}}

	\subsection{Binary Operators \showfamily}

	\begin{tabular}{*8l}
	\X\pm           &\X\cap         &\X\diamond             &\X\oplus     \\
	\X\mp           &\X\cup         &\X\bigtriangleup       &\X\ominus    \\
	\X\times        &\X\uplus       &\X\bigtriangledown     &\X\otimes    \\
	\X\div          &\X\sqcap       &\X\triangleleft        &\X\oslash    \\
	\X\ast          &\X\sqcup       &\X\triangleright       &\X\odot      \\
	\X\star         &\X\vee         &\X\lhd                 &\X\bigcirc   \\
	\X\circ         &\X\wedge       &\X\rhd                 &\X\dagger    \\
	\X\bullet       &\X\setminus    &\X\unlhd               &\X\ddagger   \\
	\X\cdot         &\X\wr          &\X\unrhd               &\X\S         \\
	\X+             &\X-            &\X\amalg               &\X\P
	\end{tabular}

	\subsection{Relations \showfamily}

	\begin{tabular}{*8l}
	\X\leq          &\X\geq         &\X\equiv       &\X\models      \\
	\X\prec         &\X\succ        &\X\sim         &\X\perp        \\
	\X\preceq       &\X\succeq      &\X\simeq       &\X\mid         \\
	\X\ll           &\X\gg          &\X\asymp       &\X\parallel    \\
	\X\subset       &\X\supset      &\X\approx      &\X\bowtie      \\
	\X\subseteq     &\X\supseteq    &\X\cong        &\X\Join        \\
	\X\sqsubset     &\X\sqsupset    &\X\neq         &\X\smile       \\
	\X\sqsubseteq   &\X\sqsupseteq  &\X\doteq       &\X\frown       \\
	\X\in           &\X\ni          &\X\propto      &\X=            \\
	\X\vdash        &\X\dashv       &\X<            &\X>            \\
	\X:
	\end{tabular}

	\subsection{Punctuation \showfamily}

	\begin{tabular}{*{5}{lp{3.2em}}}
	\X,     &\X;    &\X\colon       &\X\ldotp       &\X\cdotp
	\end{tabular}

	\subsection{Arrows \showfamily}

	\begin{tabular}{*6l}
	\X\leftarrow            &\X\longleftarrow       &\X\uparrow     \\
	\X\Leftarrow            &\X\Longleftarrow       &\X\Uparrow     \\
	\X\rightarrow           &\X\longrightarrow      &\X\downarrow   \\
	\X\Rightarrow           &\X\Longrightarrow      &\X\Downarrow   \\
	\X\leftrightarrow       &\X\longleftrightarrow  &\X\updownarrow \\
	\X\Leftrightarrow       &\X\Longleftrightarrow  &\X\Updownarrow \\
	\X\mapsto               &\X\longmapsto          &\X\nearrow     \\
	\X\hookleftarrow        &\X\hookrightarrow      &\X\searrow     \\
	\X\leftharpoonup        &\X\rightharpoonup      &\X\swarrow     \\
	\X\leftharpoondown      &\X\rightharpoondown    &\X\nwarrow     \\
	\X\rightleftharpoons    &\X\leadsto
	\end{tabular}

	\subsection{Miscellaneous Symbols \showfamily}

	\begin{tabular}{*8l}
	\X\ldots        &\X\cdots       &\X\vdots       &\X\ddots       \\
	\X\aleph        &\X\prime       &\X\forall      &\X\infty       \\
	\X\hbar         &\X\emptyset    &\X\exists      &\X\Box         \\
	\X\imath        &\X\nabla       &\X\neg         &\X\Diamond     \\
	\X\jmath        &\X\surd        &\X\flat        &\X\triangle    \\
	\X\ell          &\X\top         &\X\natural     &\X\clubsuit    \\
	\X\wp           &\X\bot         &\X\sharp       &\X\diamondsuit \\
	\X\Re           &\X\|           &\X\backslash   &\X\heartsuit   \\
	\X\Im           &\X\angle       &\X\partial     &\X\spadesuit   \\
	\X\mho          &\X.            &\X|            &\X!
	\end{tabular}

	\subsection{Variable-Sized Operators \showfamily}

	\begin{tabular}{*6l}
	\X\sum          &\X\bigcap      &\X\bigodot     \\
	\X\prod         &\X\bigcup      &\X\bigotimes   \\
	\X\coprod       &\X\bigsqcup    &\X\bigoplus    \\
	\X\int          &\X\bigvee      &\X\biguplus    \\
	\X\oint         &\X\bigwedge
	\end{tabular}

	\subsection{Log-Like Operators \showfamily}

	\begin{tabular}{*8l}
	\Z\arccos &\Z\cos  &\Z\csc &\Z\exp &
	           \Z\ker    &\Z\limsup &\Z\min &\Z\sinh \\
	\Z\arcsin &\Z\cosh &\Z\deg &\Z\gcd &
	           \Z\lg     &\Z\ln     &\Z\Pr  &\Z\sup  \\
	\Z\arctan &\Z\cot  &\Z\det &\Z\hom &
	           \Z\lim    &\Z\log    &\Z\sec &\Z\tan  \\
	\Z\arg    &\Z\coth &\Z\dim &\Z\inf &
	           \Z\liminf &\Z\max    &\Z\sin &\Z\tanh
	\end{tabular}

	\subsection{Delimiters \showfamily}

	\begin{tabular}{*8l}
	\X(             &\X)            &\X\uparrow     &\X\Uparrow     \\
	\X[             &\X]            &\X\downarrow   &\X\Downarrow   \\
	\X\{            &\X\}           &\X\updownarrow &\X\Updownarrow \\
	\X\lfloor       &\X\rfloor      &\X\lceil       &\X\rceil       \\
	\X\langle       &\X\rangle      &\X/            &\X\backslash   \\
	\X|             &\X\|
	\end{tabular}

	\subsection{Large Delimiters \showfamily}

	\begin{tabular}{*8l}
	\Y\rmoustache&  \Y\lmoustache&  \Y\rgroup&      \Y\lgroup\\[5pt]
	\Y\arrowvert&   \Y\Arrowvert&   \Y\bracevert
	\end{tabular}

	\subsection{Math Mode Accents \showfamily}

	\begin{tabular}{*{10}l}
	\W\hat{a}     &\W\acute{a}  &\W\bar{a}    &\W\dot{a}    &\W\breve{a}\\
	\W\check{a}   &\W\grave{a}  &\W\vec{a}    &\W\ddot{a}   &\W\tilde{a}\\
	\end{tabular}

	\subsection{Miscellaneous Constructions \showfamily}

	\begin{tabular}{*4l}
	\W\widetilde{abc}       &\W\widehat{abc}                        \\
	\W\overleftarrow{abc}   &\W\overrightarrow{abc}                 \\
	\W\overline{abc}        &\W\underline{abc}                      \\
	\W\overbrace{abc}       &\W\underbrace{abc}                     \\[5pt]
	\W\sqrt{abc}            &$\sqrt[n]{abc}$&\verb|\sqrt[n]{abc}|   \\
	$f'$&\verb|f'|          &$\frac{abc}{xyz}$&\verb|\frac{abc}{xyz}|
	\end{tabular}

	\subsection{{\textls[40]{AMS}} Delimiters \showfamily}

	\begin{tabular}{*8l}
	\X\ulcorner&\X\urcorner&\X\llcorner&\X\lrcorner
	\end{tabular}

	\subsection{{\textls[40]{AMS}} Arrows \showfamily}

	\begin{tabular}{*8l}
	\X\dashrightarrow       &\X\dashleftarrow
	        \\ \X\leftleftarrows      &\X\leftrightarrows     \\
	\X\Lleftarrow           &\X\twoheadleftarrow
	        \\ \X\leftarrowtail       &\X\looparrowleft       \\
	\X\leftrightharpoons    &\X\curvearrowleft
	        \\ \X\circlearrowleft     &\X\Lsh                 \\
	\X\upuparrows           &\X\upharpoonleft
	        \\ \X\downharpoonleft     &\X\multimap            \\
	\X\leftrightsquigarrow  &\X\rightrightarrows
	        \\ \X\rightleftarrows     &\X\rightrightarrows    \\
	\X\rightleftarrows      &\X\twoheadrightarrow
	        \\ \X\rightarrowtail      &\X\looparrowright      \\
	\X\rightleftharpoons    &\X\curvearrowright
	        \\ \X\circlearrowright    &\X\Rsh                 \\
	\X\downdownarrows       &\X\upharpoonright
	        \\ \X\downharpoonright    &\X\rightsquigarrow
	\end{tabular}

	\subsection{{\textls[40]{AMS}} Negated Arrows \showfamily}

	\begin{tabular}{*8l}
	\X\nleftarrow   &\X\nrightarrow \\ \X\nLeftarrow  &\X\nRightarrow \\
	\X\nleftrightarrow&\X\nLeftrightarrow
	\end{tabular}

	\subsection{{\textls[40]{AMS}} Greek \showfamily}

	\begin{tabular}{*4l}
	\X\digamma      &\X\varkappa
	\end{tabular}

	\subsection{{\textls[40]{AMS}} Hebrew \showfamily}

	\begin{tabular}{*6l}
	\X\beth &\X\daleth      &\X\gimel
	\end{tabular}

	\subsection{{\textls[40]{AMS}} Miscellaneous \showfamily}

	\begin{tabular}{*8l}
	\X\hbar         &\X\hslash      \\ \X\vartriangle &\X\triangledown      \\
	\X\square       &\X\lozenge     \\ \X\circledS    &\X\angle             \\
	\X\measuredangle&\X\nexists     \\ \X\mho         &\X\Finv$^u$          \\
	\X\Game$^u$     &\X\Bbbk$^u$    \\ \X\backprime   &\X\varnothing        \\
	\X\blacktriangle&\X\blacktriangledown \\ \X\blacksquare&\X\blacklozenge  \\
	\X\bigstar      &\X\sphericalangle     \\ \X\complement  &\X\eth       \\
	\X\diagup$^u$   &\X\diagdown$^u$
	\end{tabular}

	$^u$ Not defined in {\tt amssymb.sty}, define using the
	\verb|\newsymbol| command.

	\subsection{{\textls[40]{AMS}} Binary Operators \showfamily}

	\begin{tabular}{*8l}
	\X\dotplus      &\X\smallsetminus \\ \X\Cap        &\X\Cup               \\
	\X\barwedge     &\X\veebar      \\ \X\doublebarwedge&\X\boxminus        \\
	\X\boxtimes     &\X\boxdot      \\ \X\boxplus     &\X\divideontimes     \\
	\X\ltimes       &\X\rtimes      \\ \X\leftthreetimes&\X\rightthreetimes \\
	\X\curlywedge   &\X\curlyvee    \\ \X\circleddash &\X\circledast        \\
	\X\circledcirc  &\X\centerdot   \\ \X\intercal
	\end{tabular}

	\subsection{{\textls[40]{AMS}} Relations \showfamily}

	\begin{tabular}{*2l}
	\X\leqslant    \\\X\lesssim    \\
	\X\approxeq    \\\X\lll        \\
	\X\lesseqgtr   \\\X\doteqdot   \\
	\X\fallingdotseq\\\X\backsimeq  \\
	\X\Subset      \\\X\preccurlyeq\\
	\X\precsim     \\\X\vartriangleleft\\
	\X\vDash      \\\X\smallsmile \\
	\X\bumpeq      \\\X\geqq       \\
	\X\eqslantgtr  \\\X\gtrapprox  \\
	\X\ggg         \\\X\gtreqless  \\
	\X\eqcirc      \\\X\triangleq  \\
	\X\thickapprox \\\X\Supset     \\
	\X\succcurlyeq \\\X\succsim    \\
	\X\vartriangleright\\\X\Vdash      \\
	\X\shortparallel\\\X\pitchfork  \\
	\X\blacktriangleleft \\\X\backepsilon\\
	\X\because
	\end{tabular}

	\subsection{{\textls[40]{AMS}} Negated Relations \showfamily}

	\begin{tabular}{*8l}
	\X\nless        &\X\nleq        \\ \X\nleqslant   &\X\nleqq       \\
	\X\lneq         &\X\lneqq       \\ \X\lvertneqq   &\X\lnsim       \\
	\X\lnapprox     &\X\nprec       \\ \X\npreceq     &\X\precnsim    \\
	\X\precnapprox  &\X\nsim        \\ \X\nshortmid   &\X\nmid        \\
	\X\nvdash       &\X\nvDash      \\ \X\ntriangleleft&\X\ntrianglelefteq\\
	\X\nsubseteq    &\X\subsetneq   \\ \X\varsubsetneq&\X\subsetneqq  \\
	\X\varsubsetneqq&\X\ngtr        \\ \X\ngeq        &\X\ngeqslant   \\
	\X\ngeqq        &\X\gneq        \\ \X\gneqq       &\X\gvertneqq   \\
	\X\gnsim        &\X\gnapprox    \\ \X\nsucc       &\X\nsucceq     \\
	\X\nsucceqq     &\X\succnsim    \\ \X\succnapprox &\X\ncong       \\
	\X\nshortparallel&\X\nparallel  \\ \X\nvDash      &\X\nVDash      \\
	\X\ntriangleright&\X\ntrianglerighteq \\ \X\nsupseteq&\X\nsupseteqq\\
	\X\supsetneq    &\X\varsupsetneq \\ \X\supsetneqq  &\X\varsupsetneqq
	\end{tabular}

\fi

\clearpage  % Necessary for the last page in the appendix to display the correct page number

\end{appendix}


\ifthenelse{\equal{\manuscripttype}{BA} \OR \equal{\manuscripttype}{MA}}{%
	\renewcommand{\thepage}{}
	\ifthenelse{%
		\equal{\manuscriptlanguage}{DE}%
	}{% German
		\selectlanguage{ngerman}
		\section*{Selbstständigkeitserklärung}
		% Taken from
		% https://www.econ.uni-bonn.de/examinations/de/informationen/bachelor/bachelorarbeit/dokumente/ba-merkblatt-2016-05-23.pdf
		\pdfbookmark[1]{Selbstständigkeitserklärung}{declaration}
		Ich versichere hiermit, dass ich die vorstehende \ifthenelse{\equal{\manuscripttype}{BA}}{Bachelorarbeit}{Masterarbeit}
		selbstständig verfasst und keine anderen als die angegebenen Quellen und Hilfsmittel benutzt habe, dass die vorgelegte Arbeit noch an keiner anderen Hochschule zur Prüfung vorgelegt wurde und dass sie weder ganz noch in Teilen bereits veröffentlicht wurde. Wörtliche Zitate und Stellen, die anderen Werken dem Sinn nach entnommen sind, habe ich in jedem einzelnen Fall kenntlich gemacht. \par
		\vspace{\baselineskip}
		\noindent\datengerman\printdate{\manuscriptdate} \par
	}{% English
		\selectlanguage{USenglish}
		\section*{Declaration}
		% Taken from
		% https://www.econ.uni-bonn.de/examinations/en/information/master-economics/master-thesis/documents/ma-master-thesis-style-guide-2014-06-10.pdf
		\pdfbookmark[1]{Declaration}{declaration}
		I~hereby confirm that the work presented has been performed and interpreted solely by myself except for where I~explicitly identified the contrary. I~assure that this work has not been presented in any other
		form for the fulfillment of any other degree or qualification. Ideas taken from other works in letter and in spirit are identified in every single case. \par
		\vspace{\baselineskip}
		\noindent\printdate{\manuscriptdate} \par
	}
	\vspace{3\baselineskip} \par
	\noindent\studentfullname%
}{}


